\documentclass[parskip=full]{article}

%%%%%%%%%%%%%%%%%%%%%%%%%%%%%%%%%%%%%%%%%%%%%%%%%
% Basic includes
%%%%%%%%%%%%%%%%%%%%%%%%%%%%%%%%%%%%%%%%%%%%%%%%%


\usepackage[pdftex, colorlinks=true, linkcolor=blue, citecolor=blue, urlcolor=blue, bookmarks=false]{hyperref}

\usepackage{arydshln}
\usepackage{array}
\usepackage{makecell}
\usepackage{tikz}
\usepackage{amsmath}
\usepackage{amsfonts}
\usepackage{verbatim}
\usepackage{mathrsfs}
\usepackage{amssymb}
\usepackage{amsthm}
\usepackage{bbm}
\usepackage{bm}
\usepackage[capposition=top]{floatrow}
\usepackage[dvipsnames]{xcolor}
\usepackage{ulem}
\usepackage[labelfont=bf]{caption}
\usepackage{subcaption}
\usepackage{graphicx}
\usepackage{float}
\usepackage{natbib}
\usepackage{booktabs}
\usepackage{longtable}
\usepackage{multirow}
\usepackage{rotating}
\usepackage{pdflscape}
\usepackage[title]{appendix}
\usepackage{footmisc}
\usepackage{bigdelim}
\usepackage{tabularx}
\usepackage{xspace}
\newcommand{\A}{\ensuremath{\mathcal{A}}\xspace}
\newcommand{\B}{\ensuremath{\mathcal{B}}\xspace}
\newcommand\pa[1]{\ensuremath{\left(#1\right)}}
\usepackage{epigraph}
\usepackage{setspace}
%\onehalfspacing
\setstretch{1.5}

%%%%%%%%%%%%%%%%%%%%%%%%%%%%%%%%%%%%%%%%%%%%%%%%%%%
% Microtype
%%%%%%%%%%%%%%%%%%%%%%%%%%%%%%%%%%%%%%%%%%%%%%%%%%%
\usepackage[activate={true,nocompatibility},final,tracking=true,kerning=true,spacing=true,factor=1100,stretch=10,shrink=10]{microtype}
% activate={true,nocompatibility} - activate protrusion and expansion
% final - enable microtype; use ``draft'' to disable
%tracking=true, kerning=true, spacing=true - activate these techniques
% factor=1100 - add 10% to the protrusion amount (default is 1000)
% stretch=10, shrink=10 - reduce stretchability/shrinkability (default is 20/20)

%%%%%%%%%%%%%%%%%%%%%%%%%%%%%%%%%%%%%%%%%%%%%%%%%%%
% Hyperref setup
%%%%%%%%%%%%%%%%%%%%%%%%%%%%%%%%%%%%%%%%%%%%%%%%%%%
\usepackage{url}
\makeatletter
\g@addto@macro{\UrlBreaks}{\UrlOrds}
\makeatother
\usepackage{hyperref}
\hypersetup{
  colorlinks,%
  citecolor=black,%
  filecolor=black,%
  linkcolor=blue,%
  urlcolor=black
}

%%%%%%%%%%%%%%%%%%%%%%%%%%%%%%%%%%%%%%%%%%%%%%%%%%%
% Font setup
%%%%%%%%%%%%%%%%%%%%%%%%%%%%%%%%%%%%%%%%%%%%%%%%%%%
\usepackage{mathpazo} % math & rm
%\linespread{1}        % Palatino needs more leading (space between lines)
\usepackage{parskip}
\usepackage[scaled]{helvet} % ss
\usepackage{courier} % tt
\usepackage{natbib}

\usepackage{filecontents}

\normalfont
\usepackage[T1]{fontenc}

\usepackage{titlesec}
\titleformat{\section}
{\normalfont\fontsize{12}{17}\bfseries}
{\thesection}
{1em}
{}

\titleformat{\subsection}
{\normalfont\fontsize{11}{16}\bfseries}
{\thesubsection}
{1em}
{}

%%%%%%%%%%%%%%%%%%%%%%%%%%%%%%%%%%%%%%%%%%%%%%%%%
% macros
%%%%%%%%%%%%%%%%%%%%%%%%%%%%%%%%%%%%%%%%%%%%%%%%%
\newcolumntype{C}[1]{>{\centering\let\newline\\\arraybackslash\hspace{0pt}}m{#1}}
\newcommand{\indicator}[1]{\mathbbm{1}_{\left\{ {#1} \right\} }}
%\newcommand{\resultsloc}{../code/python/output/}
\newcommand{\subfigwidth}{0.475\textwidth}
\newcommand{\subfigwidthtwo}{0.375\textwidth}
\newcommand{\figwidth}{0.6\textwidth}
\DeclareMathOperator*{\argmax}{arg\,max}

\setlength\epigraphwidth{0.9\textwidth}
\setlength\epigraphrule{0pt}
\renewcommand{\epigraphflush}{center}
\usepackage{etoolbox}
\makeatletter
\patchcmd{\epigraph}{\@epitext{#1}}{\itshape\@epitext{#1}}{}{}
\makeatother

%%%%%%%%%%%%%%%%%%%%%%%%%%%%%%%%%%%%%%%%%%%%%%%%%
% Page and paragraph settings
%%%%%%%%%%%%%%%%%%%%%%%%%%%%%%%%%%%%%%%%%%%%%%%%%
\setlength{\parskip}{0.0pc} \setlength{\parindent}{0.25in}
\setlength{\topmargin}{-3pc}
\setlength{\textheight}{9in}
\setlength{\oddsidemargin}{0pc} \setlength{\evensidemargin}{0pc}
\setlength{\textwidth}{6.5in}
\makeatletter
\@fpsep\textheight
\makeatother

\newenvironment{fignote}{\begin{quote}\footnotesize}{\end{quote}}
\def\sym#1{\ifmmode^{#1}\else\(^{#1}\)\fi}
\renewcommand{\baselinestretch}{1.5}

% Formats the notes at the bottom of tables and figures.
\newenvironment{fignote2}[1]
{
\par
\medskip
\begin{minipage}{#1}
\footnotesize
}
{
\end{minipage}
}

% Formats the notes at the bottom of tables and figures.
\newenvironment{fignote3}[1]
{
%\par
\begin{minipage}{#1}
\footnotesize
}
{
\end{minipage}
}

%%%%%%%%%%%%%%%%%%%%%%%%%%%%%%%%%%%%%%%%%%%%%%%%%
% Title page
%%%%%%%%%%%%%%%%%%%%%%%%%%%%%%%%%%%%%%%%%%%%%%%%%
%%%%%%%%%%%%%%%%%%%%%%%%%%%%%%%%%%%%%%%%%%%%%%%%%
% Title page
%%%%%%%%%%%%%%%%%%%%%%%%%%%%%%%%%%%%%%%%%%%%%%%%%

\title{\vspace{-1.8cm}The Welfare Consequences of Rent Control\footnote{We extend our gratitude to the seminar attendees at the University of Wisconsin-Madison and the University of Windsor, as well as to Amy Sun, Joseph B. Steinberg, Yahong Zhang, Ruoyu (Steven) Chen, Chenggang (Mark) Zhou, and John Baker for their valuable insights and comments. We owe a special thanks to Joseph B. Steinberg for generously providing access to his high-powered computer for this research. The data and replication code are available on GitHub at \url{www.github.com/srotberg/the_welfare_consequences_of_rent_control}.}}

\author{Shahar Rotberg\footnote{University of Toronto. Corresponding author: \url{shahar.rotberg@mail.utoronto.ca}. Website: \url{www.sites.google.com/site/shaharrotberg}} \and Lin Zhang\footnote{Canada Mortgage and Housing Corporation. Contact: \url{lzhang@cmhc-schl.gc.ca}. The opinions expressed herein do not necessarily represent the views of CMHC.}}

\pagestyle{plain}

\begin{document}

\maketitle

\setstretch{1.5}
\begin{abstract}
\onehalfspacing
\noindent Tenancy and unit controls are common rent controls limiting rent hikes. Under tenancy control, landlords can reset rents upon tenant turnover, and thus market premiums allocate controlled units, which delays rent control's benefits and may exclude low-income households from the market. Unit control provides immediate discounts and enables low-income households market access, but controlled units are allocated via lottery. We nest both policies in a general equilibrium housing model calibrated to Toronto where tenancy control is in place. In our model, tenancy control generates higher welfare than unit control, and the welfare-maximizing degree of tenancy control is 65\%. 
\\[1.5ex]
% words=98

\noindent \textbf{JEL Classification:} E17; R31; R52. 

\noindent \textbf{Keywords:} General equilibrium, housing markets, rent control,  tenancy control, unit control, homelessness.
\end{abstract}
\setstretch{1.5}

\section{Introduction}

While tenancy control and unit control are rent control policies that are implemented across various cities globally, the economic literature has predominantly scrutinized unit control, leaving tenancy control largely unexplored. Unit control imposes limits on rent hikes for a rent-controlled unit throughout a tenancy as well as when tenants change, as observed in cities like New York City, Vancouver, Montréal, and Paris. In contrast, tenancy control regulates rent increases during a tenancy but grants landlords the freedom to set rents at their discretion upon tenant turnover. This policy is seen in cities such as Toronto, San Francisco, Los Angeles, Washington D.C., Singapore, and Hong Kong, among others.

In this paper, we delve into essential distinctions between these two categories of rent control and construct a general equilibrium housing model that nests both tenancy control and unit control and encapsulates these disparities. Our starting point is a model featuring tenancy control, calibrated using data from Toronto, where this policy is enacted. To the best of our knowledge, we offer the inaugural theoretical and quantitative analysis of tenancy control within a general equilibrium framework and are the first to compare it with unit control.

Using our model, we explore whether switching from our baseline model, which employs tenancy control, to unit control would be welfare-enhancing. Additionally, we analyze the impact on equilibrium prices and welfare within our baseline model when either fully decontrolling or fully controlling the rental market. We also search for the welfare-maximizing degree of tenancy control. Our analysis reveals that tenancy control yields higher welfare in our model, and controlling 65\% of the rental market, which closely aligns with the benchmark degree of tenancy control, maximizes welfare. Below, we discuss our findings in detail.

What are the fundamental disparities between tenancy control and unit control? First, they differ in their approaches to allocating controlled units. Unit control offers new tenants immediate access to discounted rents compared to free market rents. However, this creates a shortage of rental housing because the quantity demanded surpasses the quantity supplied at the government-mandated below-market rents. Consequently, as \citet{gl-2003} estimate, rental units are assigned arbitrarily, and thus do not necessarily go to households with the highest willing to pay. Conversely, tenancy control permits landlords to adjust rents when tenants change, and thus ensures that market forces allocate controlled units to households with the highest willingness to pay. This mechanism averts a shortage of controlled rental housing. However, it entails that new tenants in controlled units must pay a premium over free market rents to secure their accommodation, thereby postponing the advantages of rent control for the tenant.

The second distinction lies in the capacity of each policy to assign controlled units to households with the greatest marginal utility. Given that housing cannot be consumed in arbitrarily small quantities, certain low-income households, which derive substantial utility from housing, might find themselves excluded from the rental market, despite it being socially optimal to prioritize them for controlled units. Consequently, these households may face homelessness or be compelled to live with parents. Unit control offers immediate rent discounts, potentially eliminating barriers to entry into the controlled rental market for low-income households. On the other hand, under tenancy control, renters are required to pay a premium to secure a controlled unit, which may exclude certain low-income households with the greatest marginal utilities from accessing the controlled market.

To answer our questions, we construct a model that encompasses both types of rent control policies, with tenancy control serving as our foundational model. Our model is a general equilibrium housing framework featuring overlapping generations and two rental markets operating in parallel: an uncontrolled market (i.e., free market) and a controlled market. Within our baseline model, which features tenancy control, renters have the freedom to choose between renting in the controlled market or the free market. In the controlled market, rents for units are regulated throughout a renter's tenancy, gradually declining in real terms as the renter remains longer in the same unit. However, these units also suffer from deterioration due to landlords' lack of maintenance. Upon a tenant's departure, landlords can revert to charging market rents for new tenants. Conversely, units in the free market are unregulated; landlords maintain these units adequately, but rents do not decrease in real terms regardless of tenancy duration and are susceptible to stochastic fluctuations. Renters are unable to insure themselves against the financial risks associated with renting in the free market.

To capture the barriers hindering low-income households from entering into rental markets, we introduce a minimum rental size, and account for the possibility of households facing homelessness or residing with their parents.\footnote{According to estimates by \citet{spatial_distance-2020}, approximately 75\% of young adults (aged 26-34) in the U.S. reside within a 30-mile radius (approximately 48 kilometers) of at least one of their parents.} Building upon \citet{az-2024}, households experiencing homelessness do not pay rent but receive minimal housing services, akin to shelter assistance local governments typically provide. Similarly, households cohabiting with their parents reside rent-free, albeit with reduced housing consumption compared to what they could secure in the rental market.

Our framework yields three primary equilibrium implications in our baseline model with tenancy control. Firstly, to capitalize on the gradual reduction in rents over time and the insurance against rent fluctuations offered by rent control, tenants must initially pay a premium to rent in the controlled market. Secondly, as the proportion of the controlled market diminishes, intensified competition for rent-controlled units pushes the premium upward. Thirdly, high-income households, equipped to handle the financial risks associated with the free rental market and anticipating to become homeowners quickly, often opt for the free market to evade paying the premium. Conversely, low-income households, with lower risk tolerance and an expectation of prolonged rental tenure, prefer to pay the premium to enjoy the advantages of rent control.

The remainder of our model encompasses fairly standard components. Households make decisions regarding renting, homeownership, or, as previously discussed, homelessness or residing with their parents. Homeowners can access long-term mortgage contracts. Working-age households earn labor income by supplying labor inelastically and receive pensions upon retirement. Households accrue capital income from their financial assets, and labor and capital incomes subject to progressive taxation. Similar to the framework outlined in \citet{garrigaeetal-2009}, a competitive rental company, operating with convex marginal costs, supplies rental housing. Drawing inspiration from \citet{fallis-1985}, we model the rental housing supply after the structure observed in Toronto's rental market, where older construction is subject to rent control while newer construction is not. In this configuration, the free market rent represents the rental company's marginal revenue, thereby determining the quantity of rental housing supplied for profit maximization. See Section \ref{sec:rental_company} for more details on the rental company. Additionally, mirroring the approach outlined in \citet{rotberg-2022}, a competitive housing construction sector, characterized by convex marginal costs, contributes to the creation of new housing units, which are subsequently added to the depreciated housing stock.

The model incorporates both federal and local government entities. The federal government is responsible for collecting consumption and income taxes, and utilizing the generated revenue to fund pensions for retired households. Conversely, the local government collects property taxes and allocates the proceeds towards providing homeless shelters, as in \citet{abramson-2024}, and a public good that does not directly impact household utility. Given that rent control typically falls under the jurisdiction of local governments, we require budget balance only at the level of the local government, which we achieve by adjusting the property tax rate.

We calibrate the long-run equilibrium of our model using data from Toronto, Canada, where, in 2016, a portion of the rental market was subject to tenancy control while the remainder of the rental market remained free. Much of our calibration process follows standard practices; we select parameter values to align with various statistics such as the homeownership rate, and the proportion of total renters' income allocated to rent, among others. However, certain aspects of our calibration are unique to our study. Firstly, we target the percentage of renters in the controlled market, which stands at 71\% according to the 2016 Census. This figure serves as a crucial determinant for the premium that new tenants in the controlled market must pay to secure a unit in equilibrium. Secondly, we consider the proportion of renters spending over 50\% of their income on rent, also drawn from the 2016 Census. This target, akin to the approach proposed by \citet{rs-2023}, informs the determination of the minimum rental unit size, thus reflecting the barrier faced by low-income households in entering the rental market. Lastly, we set the level of housing services obtained by households living with their parents and those experiencing homelessness to correspond with the percentage of 26-34-year-old households cohabiting with parents and the percentage of homeless households, respectively. We derive the former from the 2016 Census, while we source the latter from 'The Homelessness Hub'.

In our exercises, we contrast our calibrated baseline model featuring two rental markets operating in parallel (our benchmark) with alternative long-run equilibrium scenarios. Drawing from the principles outlined in \citet{ralws-1971}, we define welfare as the ex-ante expected utility of newborns, a widely accepted convention in the literature \citep{conesa-et-al,gkkoc}.

In our initial experiment, we juxtapose our baseline model's benchmark against unit control. Under unit control in our model, households are ``wake up'' each period with an independent and random opportunity to rent in either the controlled market or the free market. Those opting for the controlled market can retain their unit indefinitely. Conversely, households lacking this opportunity may choose to rent in the free market, pursue homeownership, reside with parents if they can, or face homelessness. To ensure comparability with our baseline, we implement three adjustments. Firstly, we establish the controlled rent in the unit control scenario, mandated exogenously by the government, to mirror the average controlled rent in the baseline's benchmark. Secondly, we maintain consistency by enforcing the same proportion of the rental market to be under control in the unit control setup. Thirdly, we align the rent reduction in the unit control setup with the rate of unit deterioration resulting from landlord neglect. This approach ensures that controlled units deteriorate at an equivalent pace across both scenarios while maintaining constant real controlled rent in our unit control experiment.

Because unit control would result in a shortage of controlled units, these units would be allocated based on chance. Consequently, fewer controlled units would be allocated to low-income households compared to our benchmark with tenancy control. This circumstance would compel more high-income households to remain in the rental market rather than transitioning to homeownership, thereby increasing total rental demand and prompting a 2.4\% escalation in free market rents. However, unit control would offer an immediate rent discount to tenants in controlled units; a benefit that takes approximately 10.5 years for a tenant under tenancy control to attain. As a result, newborns' welfare would only decline by 0.06\%. Furthermore, we demonstrate that welfare gains could materialize if free market rents remained the same as in our benchmark. This suggests that the advantages associated with immediate rent discounts would outweigh the drawbacks linked to the suboptimal allocation of controlled units, were it not for the general equilibrium effects driving up free market rents for unlucky renters. We offer additional intuition using Figure \ref{fig:diagram_unit_control}.

Our second experiment compares the baseline benchmark with the long-run equilibrium wherein the government fully decontrols the rental market. Full decontrol would precipitate two primary effects. Firstly, it would diminish the appeal of renting, thereby increasing homeownership, driving up house prices, and reducing rental demand --- a consequence stemming from the extensive margin. Conversely, full decontrol would reduce the initial rent for new tenants previously renting in the controlled market. This lowered initial rent would spur an increase in total rental demand from those tenants remaining in the rental market --- a repercussion arising from the intensive margin. Overall, the impact from the intensive margin would marginally outweigh that from the extensive margin, resulting in a slight long-run rise in free market rents. The combination of eliminating the tenancy control option for certain renters, the marginally higher free market rents, and elevated house prices would collectively lower welfare. We provide further visual intuition in Figure \ref{fig:diagram_full_decontrol}.

Our third experiment compares the baseline model's benchmark with the long-run equilibrium in which all rental units are under tenancy control. Imposing tenancy control on the full market also would have extensive and intensive margin effects, but in reverse relative to the second experiment. At the extensive margin, more households would choose to rent, which would lower the homeownership rate and house prices, and thus increase rental demand. At the intensive margin, on the other hand, the higher initial rent renters, who previously rented in the free market, would have to pay would reduce the amount of rental housing they demand. Although the effect at the intensive margin would again outweigh the effect at the extensive margin, total rental demand for controlled units would be higher than total rental demand for controlled units in the benchmark. This would raise controlled rents new tenants would pay and the average controlled rent. The higher initial rent all new renters would have to pay would lead to welfare losses. Figure \ref{fig:diagram_full_control} provides further visual intuition.

Our third experiment contrasts the baseline model's benchmark with the long-run equilibrium in which all rental units fall under tenancy control. Imposing tenancy control across the entire market would also engender extensive and intensive margin effects, but in reverse relative to the second experiment. At the extensive margin, an increased number of high-income households would opt for renting, consequently lowering homeownership and house prices but raising rental demand. Conversely, at the intensive margin, the higher initial rent imposed on renters who previously rented in the free market would curtail their demand for rental housing. Although the effect at the intensive margin would once again outweigh that at the extensive margin, aggregate rental demand for controlled units would surpass that in the benchmark. Consequently, this would increase the average controlled rent, even though initial rents would remain roughly unchanged. Overall, while the decrease in house prices would boost welfare by facilitating newborns' entry into homeownership, the elimination of the option to rent in the free market would be more detrimental, leading to welfare losses. Figure \ref{fig:diagram_full_control} provides further visual intuition.

In our fourth experiment, we endeavor to identify the degree of tenancy control that maximizes welfare in our baseline model. We vary the proportion of rental housing subjected to tenancy control and assess welfare in the long-run equilibrium of each scenario. Our analysis reveals that, in our model, the welfare-maximizing degree of tenancy control stands at 65\%, a figure remarkably close to our benchmark where approximately 70\% of the rental market falls under tenancy control. Low degrees of tenancy control prove detrimental to welfare as they exacerbate the scarcity of controlled units, driving up their premiums and further excluding low-income households from the market. Conversely, high degrees of tenancy control also reduce welfare. Beyond a certain threshold, compelling more households to rent in the controlled market raises total rental demand, which subsequently increases free market rents and adversely impacts households preferring to rent outside the controlled sector. Lastly, consistent with the findings of our first experiment, we ascertain that unit control would diminish welfare under the same degree of control and the same average controlled rents as in our welfare-maximizing long-run equilibrium.

This paper makes significant contributions both theoretically and quantitatively to the field of rent control research. Theoretically, we uncover fundamental economic distinctions between tenancy control and unit control, which, to our knowledge, have not been previously identified or discussed in the literature. Additionally, we develop a general equilibrium model that includes both tenancy control and unit control, and thus provide a versatile framework for analysis. Quantitatively, we leverage our model to assess the welfare implications of these two types of rent control, shedding new light on their comparative effectiveness. Furthermore, we pinpoint the welfare-maximizing degree of tenancy control in our model, and thus offer valuable insights into policy optimization in this domain.

\section{Related literature}\label{sec:related_lit}

The classical theory in economics, discussed in \citet{suen-1989}, posits that government-imposed rent ceilings (termed unit control in our study) cause demand to surpass supply, and thus induce a shortage of controlled units. \citet{gl-2003} contend that when such shortages arise, the allocation of controlled units through a lottery system leads to a misallocation. However, we argue in our paper that this classical theory does not hold true under tenancy control. With tenancy control, also known as vacancy decontrol or third generation rent control, landlords can adjust rents freely when new tenants move in. Consequently, competition among potential tenants for these units raises their price, ensuring that only those willing to pay the most secure a controlled unit. Thus, under tenancy control, market forces, rather than a lottery, dictate the allocation of controlled units, resulting in a more efficient allocation and no shortage of controlled rental housing.

As highlighted in \citet{mm-2023}, the scarcity of empirical research on rent control arises from two main factors: "First, micro-level data on house prices and rents is mostly available for recent periods, whereas sources of potentially exogenous variation are available for earlier episodes. Second, there are various ways in which rental cap policies are implemented, and hence, we have only a partial understanding on how certain types of rental cap policies may affect the market." Therefore, it is imperative to employ equilibrium housing models to learn about the impacts of different rent control policies. \citet{fallis-1985} examined Toronto's rent control policies, albeit with a static model rather than a dynamic one like ours. Similar to our model, theirs included an interaction between controlled and free markets. More recent models have turned to general equilibrium models with dynamic lifecycle choices. \citet{eric-1995} and \citet{fmn-2023} modeled rent control akin to unit control and assessed its effects on house prices, rents, and welfare.\footnote{Controlled units in \citet{fmn-2023} offer immediate discounts that increase with tenure duration.} However, neither of them incorporated tenancy control as we do, nor did they compare tenancy control to unit control. Moreover, although \citet{ms-1993} show that rent control leads to a deterioration in the quality of rental housing, neither \citet{eric-1995} nor \citet{fmn-2023} allow controlled units to deteriorate as tenants remain longer in their unit. 

Our work is also related to a couple of theoretical studies with a focus on modelling tenancy control. \citet{be-2000} model tenancy control in an asymmetric information setup in which landlords do not know the types of tenants or how long tenants stay. While their model is focused on strategic interactions between landlords and tenants, our work provides a theory on how controlled units are allocated when free and tenancy-controlled markets operate at the same time, and we compare it to unit control. In addition, \citet{diamond-2019} evaluate the welfare impact of expansions in tenancy control in San Francisco using a dynamic, neighborhood choice model. Their model features many desirable features, such as tenant buyouts, the potential for biased renters' expectations about the future, and endogenous landlord responses. However, it does not feature general equilibrium elements that we focus on, such as endogenous rents, feedback effects between free market rents and controlled rents, or interactions between the rental and homeownership markets.

Our theory provides a potential explanation for why several studies have observed rent premiums in controlled markets. \citet{nagy-1997} discovered that rent-controlled apartments initially commanded higher rents than free-market ones, while \citet{early-2000} and \citet{cjq-2023} estimated that tenants in the controlled market were more likely to pay rent premiums during their initial years of tenancy. Although all these studies examined rent control in New York City where unit control is in place, it suggests that in many cases, landlords may be breaking the law and actually resetting rents upon tenant turnover. When viewed from the perspective that, in practice, unit control may actually result in equilibria that are closer to tenancy control, our model simulations of full decontrol may also align with findings from empirical studies on Cambridge, Massachusetts where unit control was repealed. For instance, both \citet{sims-2007} and \citet{app-2014} observed that the decontrolling in Cambridge in 1995 led to rent increases. Moreover, our observations in the Toronto data and in our benchmark model mirror patterns consistent with other empirical literature on rent controls. Specifically, \citet{oslen-1972} and \citet{gl-1989} found that occupants of rent-controlled apartments tended to have lower incomes than occupants of free-market units, a finding we corroborate using 2016 Census data, which show that tenants in controlled units in Toronto had incomes 19\% lower than those in free-market units.

Lastly, our work also contributes to the literature on the relationship between rent control and homelessness. First, our simulation results align with the empirical findings of \citet{eo-1998}, who estimated that rent control reduces homelessness relative to the free market. Second, while \citet{az-2024} and \citet{abramson-2024} have developed general equilibrium models of homelessness, to the best of our knowledge, we are the first to integrate homelessness into a general equilibrium model with tenancy and unit controls, and investigate the connection between the types of rent control and the homelessness rate.

%This paper contributes both theoretically and quantitatively to the research on rent control. Theoretically, we identify fundamental economic differences between tenancy control and unit control, which, to the best of our knowledge, have not been identified or discussed in the literature before.  According to the classic theory of shortages \citep{suen-1989}, rent control leads to housing shortage because demand exceeds supply, and that rationing housing may not be allocated to the tenants who are willing to pay most for them or by the price mechanism but by queues or lotteries. This is a major depart from the welfare analysis of rent controls in introductory economics, largely influenced by the studies of rent freezes during wartime or the first generation of rent controls such as \citet{hayek-1929} and \citet{fs-1946}. We argue that housing shortage occurs because of unit rent control or the second generation of rent controls \citep{arnott-1995}, which is the type of rent control and the resulted housing misallocation that \citet{gl-2003} studied in New York City. However, under tenancy control or the third generation of rent controls, vacant units are not subject to rent control and allocated to the tenants with the highest willingness to pay for them. Therefore, vacant units under tenancy control are allocated by the price mechanism rather than lotteries. 

%Quantitatively, we build a general equilibrium model that nests both tenancy control and unit control, and compare the equilibrium effects of these rent controls. Since most rent controls were adopted in the 1970s as a response to rapid inflation, it is a challenge to establish the causal effects of rent controls due to limited data \citep{mm-2023}. Moreover, it is crucial to capture general-equilibrium effects of rent controls but it is a challenge for empirical models. Using general equilibrium models can overcome these challenges. \citet{fmn-2023} evaluated the effect of housing policies such as housing vouchers and rent control on the household well-being in a dynamic stochastic spatial equilibrium model, calibrated to the New York metropolitan statistical area. The model economy includes rent stabilization and free market units, and controlled units are allocated by lotteries to tenants. \citet{eric-1995} developed a lifecycle general equilibrium model of the economy with unit rent control and no rent exemptions. 

%In our model, rent exemptions are allowed as in \citet{fallis-1985} that used a static demand and supply model with controlled and free markets interacting with each other; housing quality varies as in \citet{frankena-1975} that proposed a model of rent control with variable quality. In our model, tenancy-controlled units are allocated by the price mechanism and unit-controlled ones by lotteries. Most of all, \citet{fmn-2023} is focused on policy experiments of several housing policies, while we focus on modeling tenancy control, comparing the welfare implications of the two types of rent control, and finding the welfare-maximizing degree of rent control under tenancy control. 

%Our work is related to a couple of theoretical studies with a focus on modelling tenancy control.  \citet{be-2000} model tenancy rent control in an asymmetric information setup where landlords do not know the types of tenants or how long tenants stay. They show that when tenants’ types are exogenous tenancy rent control results in a Pareto sub-optimal equilibrium, whereas free market is Pareto optimal. When the types of tenants are endogenous, tenancy rent control may give rise to multiple equilibria of low or high rents. While their model is focused on strategical interactions between landlords and tenants, our work provides a theory on how controlled units are allocated when free and tenancy-controlled markets operate at the same time, and we compare it to unit control. In addition, \citet{diamond-2019} evaluate the welfare impact of expansions in rent controls in San Francisco using a dynamic, neighborhood choice model. Their model features many desirable features, such as tenant buyouts, the potential for biased renters' expectations about the future, and endogenous landlord responses. However, it does not feature general equilibrium elements that we focus on, such as endogenous rents, feedback effects between free market rents and controlled rents, or interactions between the rental and homeownership markets. As \citet{diamond-2019}, we also find rent controls expansions would raise average rents and hurt newborns.

%Our work is also related to the literature on the relation between rent control and homelessness. Regarding the relation between rent control and homelessness, there are two effects in place. On the one hand, rent control leads to a lower vacancy rate and higher rental prices in the uncontrolled sector, which leads to more homelessness. On the other hand, rent control reduces rental prices in the controlled sector and thereby reduces homelessness \textcolor{red}{SR: Im not sure if these tradeoffs are exactly correct. It really depends on the kind of policy.}. \textcolor{red}{Shahar: could you add the two GE papers on homelessness here? Lin: they are these ones: \citet{az-2024} and \citet{abramson-2024}}.  Empirically, \citet{eo-1998} uses microdata in the US on households to estimate the effect of eliminating rent control on the probability of homelessness. Their results suggest that rent control reduces homelessness.  

%Our work also predicts the equilibrium effects of rent decontrol, evaluated in several empirical studies. \citet{sims-2007} finds that rent decontrol in Cambridge, Massachusetts in 1995 spurred rent increases, and raised the quality and quantity of rental housing;  \citet{app-2014} find that the same rent decontrol generated substantial price appreciation at decontrolled and never-controlled units, and that the majority of price appreciation came from never-controlled units or spill over effect. We find that older households and households that have already lived in controlled units benefit more from rent controls which are also observed in \citet{dmq-2019} that evaluate the impacts of tenancy control in San Francisco. 

\section{Institutional background on rent control in Toronto, Canada}\label{sec:insitutional_background}



The National Housing Act introduced rent controls in Canada in 1944 during World War II. This “first-generation” of rent control froze nominal rents and was inspired by the rent freeze in European countries during World War I intended to protect veterans from evictions and rent increases. These rent controls were imposed during wartime, when a rapid increase in demand for rental units was associated with prohibitions against building any new units to save materials for war production\citet{downs-1988}. In the same spirit, rent freezes were imposed during COVID-19. The rent freeze seemed justified during the war, but many economists believed that maintaining it after had significant negative consequences, including reduced rental supply, labor mobility, and property maintenance \citep{arnott-1995}. As a result, the Ontario provincial government repealed the “first-generation” of rent controls in 1949.  

In 1975, during the oil crisis, the Ontario government re-introduced rent controls as part of the Residential Premises Rent Review Act. These new controls, classified as “second-generation” rent controls, differed in two ways from “first-generation” rent controls. First, they allowed mild nominal rent increases once a year based on mandated guidelines. Moreover, courts could also grant increases beyond the statutory guidelines in cases of major increases in the cost of supplying rental housing, financial losses of landlords, or required capital expenditures. Second, they exempted luxury accommodation, government owned buildings, social housing, and private rental units built after January 1976.

Since 1975, the Ontario rent control system has been evolving with two major changes. The first is intermittent exemptions. In 1985, the Residential Rent Regulation Act removed exemptions for previously exempted units and newly built units. Under this new change, both new buildings and buildings constructed after 1976 were subject to rent controls. In 1992, rent control with exemptions for new construction was again restored by The Rent Control Act, which exempted all rental buildings built after November 1, 1991. However, in April 2017, the government of Ontario removed exemptions by extending rent controls to rental units built after November 1, 1991. In 2018, the Ontario government enacted new legislation exempting only buildings built, or first occupied, after November 15, 2018.

The second major change in the Ontario rent control system came in 1998, and was marked by tenancy deregulation as part of the Ontario Tenant Protection Act. This Act allowed landlords of rent-controlled units to increase rents by more than prescribed by the rent increase guidelines, but only when signing a new contract with a new tenant. Otherwise, legislation compelled them to follow the rent increase guidelines. This type of rent control is sometimes referred to as tenancy rent control, partial rent control, or “third-generation” rent control \citep{arnott-2003}. The Ontario government has not made any changes since. 

\textcolor{red}{review of other cities with rent control}

\section{Model}\label{sec:model}

The model encompasses overlapping generations of finitely-lived households, competitive rental and housing construction firms, and governmental bodies. Each age cohort of households exhibits heterogeneity in labor market productivity, housing, mortgage debt, financial assets, and housing tenure (ownership, renting, cohabiting with parents, or homelessness). Renters further differ in the rental market type (controlled or free), along with the duration of their residence in a particular unit. Renters in the free market face rent shocks.

Pre-retirement, households engage in the labor market, while post-retirement, they receive government pensions dependent on their average lifetime labor income. Renters make decisions regarding rental market selection and tenure duration. Homeowners incur proportional property taxes and maintenance expenses, while we treat mortgages as long-term contracts, following the approach of \citet{kkk-2021}. Until age $J_{p}$, households may opt for rent-free accommodation with their parents, albeit at the cost of reduced housing services compared to the rental market. Past age $J_{p}$, households ineligible for parental cohabitation face homelessness if they are unable or unwilling to engage in the rental market, receiving minimal housing services at no cost.

Competitive rental firms, characterized by convex marginal costs, operate in both controlled and free rental markets simultaneously. Similarly, competitive construction companies, subject to increasing marginal costs, invest in new housing production, which supplements the depreciated housing stock.

The model integrates both federal and local government entities. The federal government finances pensions for retirees through income and consumption taxes, whereas local governments offer services to residents that they fund by imposing property taxes. Given the localized nature of rent control legislation, we require budget balance solely at the local government level.

Considering that our model economy is a city, we assume that the interest rate, $r$, is exogenous \citep{sommeretal-2013, arslanetal-2015, sommersullivan-2017, kaplanetal-2019}. Household savings encompass bonds and stock investments in rental and construction firms, with an implicit requirement of indifference between these options due to a no-arbitrage condition. Because asset market clearance is not required, we do not explicitly track dividends associated with stock holdings. Alternatively, one can think of the rental and construction firms as foreign-owned, as posited by \citet{floettotetal-2016}

\subsection{Demographics}\label{subsection:demographics}

Households enter the model at age $j=1$, retire from the labor force at age $J_{R}$, and live for $J$ periods. Until age $J$, they face no mortality risk, but die with certainty afterward. Household preferences are characterized by the following utility function:

\begin{equation*}
    u_{j}(c,h(o,f,\iota))=\xi_{j}^{\sigma}\frac{\left(c^{1-\gamma} \left(h(o,f,\iota) \right)^{\gamma}\right)^{1-\sigma}}{1-\sigma},
\end{equation*}

\noindent where $c$ denotes non-housing consumption, $h$ represents housing services consumption, and $f=1$ indicates households renting in the free market, while $f=0$ denotes those renting in the controlled market.\footnote{For homeowners, households living with their parents, or homeless households, $f$ holds no significance, and thus we set it to $1$.} The parameter $\iota$ signifies the number of years a renter in the controlled market has resided in the same unit, whereas $o=1$ signifies homeowners and $o=0$ signifies renters, households cohabiting with their parents, and homeless households. The weight $\gamma$ denotes the utility assigned to shelter service consumption, $\sigma$ represents the household's constant relative risk aversion, and $\xi_{j}$ scales equivalently to accommodate variations in household size over its lifecycle, drawing from the works of \citet{lm-1980}, \citet{wz-2010}, \citet{xu-2013}, and \citet{rs-2023}. Household discounting of future consumption occurs at a rate of $\beta$, while non-housing consumption incurs a proportional tax $\tau_{c}$.

Similar to \citet{cagetti-denardi-2006}, households derive a warm-glow from bequests they leave behind upon death, which we capture by $\omega_{B}(q)=\mu \frac{q^{1-\sigma}}{1-\sigma}$. Here, $q$ represents the household's net worth at the end of their life, which is further defined below, while $\mu$ modulates the strength of the bequest motive.

\subsection{Heterogeneity}

Households are heterogeneous across multiple dimensions:

\begin{itemize}
    \item Financial assets, $a\in \mathcal{A}=\mathbb{R}_{+}$.
    \item Cohabitation status with parents, $v\in \mathcal{V}={0,1}$.
    \item Homelessness status, $b \in \mathcal{B}={0,1}$.
    \item Tenure type (ownership or rental), $o \in \mathcal{O}={0,1}$, if obtaining housing services from the housing market.
    \item Rental market selection for those receiving rental services, $f \in \mathcal{F}={0,1}$.
    \item Duration of residence in the same unit for renters in the controlled market, $\iota \in \mathcal{I}=\mathbb{N}$.
    \item Rent shock in the free market for those renting, $z \in \mathbb{Z}$.
    \item House size, $h\in \mathcal{H}=\mathbb{R}_{++}$.
    \item Mortgage debt, $m \in \mathcal{M}=\mathbb{R}_{+}$.
    \item Idiosyncratic labor fixed effect $\chi \in \mathcal{X}$.
    \item Stochastic labor market productivity shock, $e \in \mathcal{E} = \mathbb{R}_{++}$.
\end{itemize}
 
\noindent We denote the measure of age-$j$ households over endowments by $\Psi_{j}(a,v,b,o,f,\iota,z,h,m,\chi,e)$. For brevity, we mainly employ $s$ to represent a household's state $(a,v,b,o,f,\iota,z,h,m,\chi,e)$.

\subsection{Income and income taxes}\label{subsection:income}

Working-age households ($j<J_{R}$) supply labor inelastically. Age-$j$ households, characterized by labor market productivities $\chi$ and $e$, earn labor income $y_{\ell,j}(\chi,e)=\chi e\zeta_j$. The fixed-effect $\chi$ remains constant throughout a household's life and is drawn from an unconditional distribution, which we denote as $\hat{F}(\chi)$. The idiosyncratic shock $e$ follows a Markov process conditional on the homelessness status, which we denote as $F(e,e'|b)$, where $b=1$ indicates homelessness and $b=0$ indicates non-homelessness. Particularly, in line with \citet{az-2024}, households experiencing homelessness are more prone to encountering persistently low stochastic productivity shocks. $\zeta_j$ represents a common life-cycle profile. Retired households receive pensions, $y_{R}(\chi,e)$, which is contingent upon labor productivity at retirement.\footnote{Retirees with higher productivity levels are more likely to have had higher productivities over their lifetimes, and thus receive larger pensions.}

Households can also earn income from their savings, $y_{k}(a)=ar$, where $r$ is the interest rate. We define taxable income as:

\begin{equation*}\label{eq:taxable_labor_income}
  y^{x}_{j}(\chi,e,a) =y_{\ell,j}(\chi,e)+\nu y_{k}(a),
\end{equation*}

\noindent where $\nu$ signifies the taxable portion of capital income. We apply progressive taxes to taxable income as follows:

\begin{equation}\label{eq:income_tax}
  y^{\tau}_{j}(\chi,e,a) = y^{x}_{j}(\chi,e,a)-\tau_{\ell}y^{x}_{j}(\chi,e,a)^{\psi},
\end{equation}

\noindent where $\tau_{\ell}$ determines the average tax rate and $\psi$ governs the degree of progressivity \citep{hsv-2020}.

As in \citet{kkk-2021}, each period, surviving households receive financial wealth endowments equivalent to the bequests left by deceased households in the preceding period. Surviving households across all ages receive financial endowments, which we denote as $\vartheta (\chi,e)$, capturing the correlation between bequests and income.

\subsection{Housing}\label{subsection:housing_choice}

Households have the option to either rent or own a dwelling, denoted by $o \in \mathcal{O}={0,1}$, cohabit with their parents ($v \in \mathcal{V}={0,1}$) if $j<J_{p}$, or face homelessness ($b \in \mathcal{B}={0,1}$) if $j \geq J_{p}$. Households that opt to rent or own, can also select the size of their dwelling, which we denote by $h$. The smallest rental unit is smaller than the smallest housing unit for purchase, defined as $0<\underline{h}_{r}<\underline{h}_{o}$ \citep{gervais,sommeretal-2013,sommersullivan-2017,kaplanetal-2019}. Households residing with their parents receive housing services equivalent to $\underline{h}_{p}$, while those choosing homelessness receive housing services equivalent to $\underline{h}_{h}$, where $0<\underline{h}_{h}<\underline{h}_{p}<\underline{h}_{r}$. Households living with their parents ($v=1$) or experiencing homelessness ($b=1$), receive their housing services free of charge. However, the city incurs a cost of $\theta_{b}$ to support each homeless household.

Households opting to own a house of size $h$ incur the purchase price $p h$. Transaction involving buying and selling houses involve one-time proportional transaction costs, $\tau_{b} p h$ and $\tau_{s} p h$, respectively. Each period, all housing units are subject to taxation at a rate $\tau_{p}$ and depreciate in value at a rate $\delta_h$.

Renters need to choose between renting in the controlled market or the free market. $f=0$ denotes renting in the controlled market, whereas $f=1$ signifies renting in the free market.

In the free market, renters leasing housing services $h$ pay $p_{r,1} h$. Rent remains unaffected by tenancy duration and is subject to uncertain shocks, $z$, following a distribution $G(z)$. With probability $G(z=x)$, renters paying $p_{r,1}h$ in the free market in the current period face additional rent payments of $x p_{r,1}h$ in the subsequent period, where $x$ represents an independent shock that can of take negative as well as positive values. Consequently, the average per-unit rent in the free market is given by $\int_\mathcal{Z}(1+z)p_{r,1}\mathrm{d}G(z)$, denoted as $\overline{p}_{r,1}$.

In contrast, rent in the controlled market is certain and diminishes with tenancy duration. Similar to \citet{fmn-2023}, prolonged tenancy in the controlled market leads to lower real rent. Following the approach of \citet{fmn-2023}, controlled rent is determined by tenancy duration as follows:

\begin{equation}\label{eq:rent_in_controlled_market}
\frac{p_{r,0}}{(1+\kappa_{1})^{\iota}} h,
\end{equation} 

\noindent where $h$ denotes the condition of the controlled unit at the time a new tenant moves in. $p_{r,0}$ represents the initial rental price paid by new tenants, while $\kappa_{1}$ signifies the rate at which rent control diminishes real rents for each additional period tenants remain in the same unit. $\iota$ indicates the number of years they have resided in that unit.

Unlike landlords of free market units who fully offset the depreciation of their properties each period, landlords of controlled units neglect their units. Consequently, controlled units deteriorate over time, resulting in a reduction of housing services available to tenants \citep{gl-1990,la-2018}. The amount of housing services tenants obtain in the controlled market is thus given by:

\begin{equation}
    h(o,f,\iota)=\frac{1}{(1+(1-o)(1-f)\kappa_{2})^{\iota}} h.
\end{equation}
 
\noindent $\kappa_{2}$ represents the rate at which controlled units deteriorate for each additional period a household remains renting the same unit.

Figure \ref{fig:diagram_rent_path} provides a qualitative illustration of rents paid by renters as a function of the market in which they rent and their tenancy duration.\footnote{To maintain clarity in our illustration, we abstract from the effects of rental unit deterioration.} The downward sloping blue line depicts rents paid by renters residing in controlled units. Upon entering a new rental unit, their rent is $p_{r,0}$, but as their tenancy duration increases, their rent decreases in accordance with equation \eqref{eq:rent_in_controlled_market}.

Conversely, the horizontal red line represents the average rents paid by renters dwelling in the free market, which we denote as $\overline{p}_{r,1}$. This average rent remains constant irrespective of their tenancy duration, although their actual rent may fluctuate around this average value.

Figure \ref{fig:diagram_rent_path} indicates that in equilibrium, the average rent paid in the free market, $\overline{p}_{r,1}$, must be lower than the initial rent paid by renters in the controlled market, $p_{r,0}$. However, the average rent in the controlled market, $\overline{p}_{r,0}$, is lower than $\overline{p}_{r,1}$. What accounts for this equilibrium outcome?

We explain this equilibrium outcome as follows: If the initial rents paid by new tenants in the controlled market were lower than the average rents paid by renters in the free market, $\overline{p}_{r,1}$, all renters would opt to rent in the controlled market. To ensure market equilibrium, renters in the controlled market must initially pay a premium. Renters who anticipate remaining in the same unit for only a short duration and have sufficient risk tolerance for rent fluctuations would choose to rent in the free market to avoid paying this premium. Conversely, renters expecting to stay in the same unit for an extended period, and whose risk tolerance for rent volatility is low, would be willing to pay the premium to benefit from protection against rent fluctuations and lower long-term rents.

\subsection{Mortgages}\label{subsection:morgtages}

In addition to rent control, the availability of mortgage contracts significantly influences households' decisions to buy or rent and is another quantitatively important economic factor. Following \citet{kkk-2021}, we model mortgages, denoted as $m$, as long-term debt instruments with associated origination costs, prepayment penalties, and borrowing restrictions. We set the mortgage interest rate to $r_{m}=r+\Delta_m$, where $r$ represents the return on financial assets, and $\Delta_m$ signifies the mortgage debt interest rate premium \citep{sommeretal-2013,sommersullivan-2017,kaplanetal-2019}. Households cannot deduct mortgage interest payments from their taxable income because mortgage interest payments are not tax-deductible in Canada.

To obtain a new mortgage, households must incur a fixed cost denoted as $\omega_{1}$. At the time of a new mortgage origination only, the loan-to-value ratio (LTV) must not exceed a fraction $\lambda_{m,1}(ph)<1$, which depends on the value of the house at the time of purchase, denoted as $ph$:

\begin{equation}\label{eq:LTV}
\frac{m}{p h} \leq \lambda_{m,1}(ph).
\end{equation}

\noindent A constraint, which we refer to as the payment-to-income ratio (PTI), also imposes limitations on mortgage debt at origination:

\begin{equation}\label{eq:PTI}
\frac{\tau_{p}ph+(r_{m}+\upsilon_{1})m}{y_{\ell,j}(\chi,e)+y_{k}(a)} \leq \lambda_{m,2}.
\end{equation}

\noindent $\upsilon_{1}$ represents the fraction of debt paid every year on average over the life of a mortgage, while $\lambda_{m,2}$ denotes the maximum PTI ratio permitted. Constraint \eqref{eq:PTI} stipulates that property taxes, interest payments, and the average fraction of the principal paid cannot exceed a fraction $\lambda_{m,2}$ of their pre-tax labor and capital incomes.

Each period, a mortgage holder must allocate at least a fraction $\upsilon_{1}$ of their total mortgage debt towards repayment and is permitted to pay up to a fraction $\upsilon_{2}>\upsilon_{1}$ of the debt without incurring a penalty. However, any payment exceeding $\upsilon_{2}m$ is subject to a prepayment penalty, which we outline as follows:

\begin{equation}\label{eq:prepayment_penalty}
\tau_{pp}(m,m')=\omega_{2}\mbox{max}\left\{(1-\upsilon_{2})m-m',0\right\}.
\end{equation}

\noindent If a mortgage holder pays less than $\upsilon_{1}m$, we consider this akin to initiating a new mortgage. Consequently, they must incur the origination cost, $\omega_{1}$, and adhere to constraints \eqref{eq:LTV}--\eqref{eq:PTI} once again.

\subsection{Dynamic program}\label{subsection:household_optimization}
Given its state variables, a household must make decisions regarding various aspects such as saving ($a'$), homeownership or renting ($o' \in \left\{0,1\right\}$), cohabitation with their parents ($v' \in \left\{0,1\right\}$), or homelessness ($b \in \left\{0,1\right\}$). Both homeowners and renters determine their housing consumption ($h'$). Renters additionally decide whether to rent in the controlled or free market ($f' \in \left\{0,1\right\}$, where we set $f'=1$ if $o'=1$, $v'=1$, or $b'=1$). Renters renting in the controlled market also choose whether to extend their tenancy by one period ($\iota'$). Homeowners decide on their mortgage debt level ($m'$) and whether to initiate a new mortgage ($n$). We formulate the dynamic program encapsulating these decisions as follows:

\begin{align}
  \label{eq:vf}
  V_{j}(s)=\max_{\{c,a',v',b',o',f',\iota',h',n,,m',q'\}}\ \biggl\{u_{j}\left(c,h'\right(o',f',\iota'\left)\right)+\indicator{j=J}\omega_{B}(q')+
  \indicator{j<J}\beta \int_{\mathcal{E}\times \mathcal{Z}}V_{j+1}(s')dF(e,e'|b')dG(z)\biggr\}.
\end{align}
Subject to
\begin{flalign}
  & c(1+\tau_{c})+a'+or_{m}m+(1-b')(1-v')(1-o')\left[(1-f')\frac{p_{r,0}}{(1+\kappa_{1})^{\iota'-1}}h'(0)+f'p_{r,1}h'\right]+\label{eq:bc} \\ & \quad (1-b)(1-v)(1-o)fzp_{r,1}h+o\left(\tau_{p}+\delta_{h}+\indicator{o'=0 \vee h'\neq h}\tau_{s}\right)ph + o'(1+\tau_{b}\indicator{o=0 \vee h'\neq h})ph'+ \notag \\ & \quad\omega_{1}n+(1-n)\tau_{pp}(m,m') + \indicator{o'=1 \land j=J}ph(\tau_{p}-\delta_{h}-\tau_{s})= \notag \\ & \quad y_{\ell,j}(\chi,e)+y_{k}(a)-y^{\tau}_{j}(\chi,e,a)+a+o(ph-m)+o'm'+\vartheta(\chi,e) \notag \\
  & a' \geq 0 \label{eq:financial_wealth} \\
  & h' \geq (1-b')(1-v') o' \underline{h}_{o}+(1-b')(1-v')(1-o')\underline{h}_{r} \label{eq:min_housing}\\
  & h' = b'\underline{h}_{h}+v'\underline{h}_{p} \label{eq:housing_parents_homeless}\\
  & v'=0 \mbox{ if } j\geq J_{p} \label{eq:can_live_with_parents}\\
  & b'=0 \mbox{ if } j<J_{p} \label{eq:can_be_homeless}\\
  & h'(o',f',\iota')=\frac{1}{(1+(1-o')(1-f')\kappa_{2})^{\iota'-1}}h' \label{eq:housing servies_controlled}\\
  & n \in \left\{0,o'\indicator{j<J_{R}}\right\} \label{eq:origination} \\
  & nm' \leq \lambda_{m,1}(ph') p h' \label{eq:LTV_1} \\
  & n (\tau_{p}ph'+(r_{m}+\upsilon_{1})m')\leq \lambda_{m,2}(y_{\ell,j}(\chi,e)+y_{k}(a))\label{eq:PTI_1} \\
  & (1-n)m'\leq (1-\upsilon_{1})m \label{eq:settle_debt} \\
  & q' = a' + o' \left[p h'-m'\right] \label{eq:networth}
\end{flalign}

\noindent Household utility includes three components: a flow utility from housing and non-housing consumption; a warm-glow flow from bequests left behind by households who die; and future value is discounted by $\beta$.\footnote{We include labor market productivity shocks for retirees as well as working-age households in our formulation of the dynamic program for notational brevity. Writing the dynamic program this way allows us to use a single equation to represent both retired and working-age households' problems.} Rentals units rented by households in the controlled market deteriorate over time due to lack of maintenance by landlords, and this deterioration shows up as a reduction in the consumption of housing services.

Household utility comprises three elements: a flow utility derived from housing and non-housing consumption, a warm-glow component stemming from bequests left behind by deceased households, and a discounted future value.\footnote{We formulate the dynamic program in a manner that accommodates labor market productivity shocks for retirees and working-age households, ensuring brevity in notation. This approach enables us to represent both retired and working-age households' using a single equation.}

Rentals units occupied by households in the controlled market deteriorate over time due to landlords' neglect to maintain them, resulting in a decline in the consumption of housing services.

Equation \eqref{eq:bc} delineates the household's budget constraint. The left-hand side encompasses the following expenditures:
\begin{itemize}
    \item Consumption and consumption taxes, $c(1-\tau_{c})$.
    \item Saving, $a'$.
    \item Interest payments for mortgage holders, $r_{m}m$.
    \item Cost of renting a unit in the controlled market for renters in that market, $\frac{p_{r,0}}{(1+\kappa_{1})^{\iota'-1}}h'$, where $h'$ represents the condition of the controlled unit upon the renter's move-in time.
    \item Cost of renting in the free market for renters in that market, $p_{r,1}h'$.
    \item Rent shocks, $z p_{r,0}h$, for households who rented in the free market in the previous period.
    \item For homeowners, property taxes, depreciation costs, and selling costs if they sell their residence, $(\tau_{p}+\delta_{h}+\indicator{o'=0\vee h'\neq h}\tau_{s})ph$.
    \item Home purchasing costs for homebuyers, $(1+\tau_{b}\indicator{o=0\vee h'\neq h})ph'$.
    \item Origination costs for homeowners initiating a new mortgage, $\omega_{1}$.
    \item Prepayment penalty costs for mortgage holders paying more than $\upsilon_{2}$ of their mortgage, $\tau_{pp}(m,m')$.
    \item Property taxes, depreciation, and selling costs for households who reach age $J$ and opt to end their lives as homeowners.
\end{itemize}

On the other hand, the right-hand part of the budget constraint comprises the following resources available for expenditure:
\begin{itemize}
    \item After-tax labor and capital incomes, $y_{\ell,j}(\chi,e)+y_{k}(a)-y^{\tau}_{j}(\chi,e,a)$.
    \item Financial assets, $a$.
    \item Housing equity for owners, $(ph-m)$.
    \item Mortgage funds raised, $m'$, for mortgage holders.
    \item Bequests $\vartheta(\chi,e)$ contingent upon the household's labor productivity.
\end{itemize} 

Constraint \eqref{eq:financial_wealth} stipulates that households lack access to unsecured debt. Constraint \eqref{eq:min_housing} mandates that all owners must consume at least $\underline{h}_{o}$ units of housing, while all renters must consume at least $\underline{h}_{r}$ units. Constraint \eqref{eq:housing_parents_homeless} says that homeless households must consume $\underline{h}_{h}$ units of housing, and households cohabiting with their parents must consume $\underline{h}_{p}$ units of housing. Constraint \eqref{eq:can_live_with_parents} specifies that households whose age is greater than or equal to $J_{p}$ cannot reside with their parents. Constraint \eqref{eq:can_be_homeless} states that households whose age is strictly less than $J_{p}$ can live with their parents and therefore cannot be homeless.

Constraint \eqref{eq:housing servies_controlled} says that renters in the controlled market experience reduced housing services consumption as they extend their stay in the same unit. Constraint \eqref{eq:origination} asserts that a household can only initiate a new mortgage if it is of working age. Constraints \eqref{eq:LTV_1} and \eqref{eq:PTI_1} represent the LTV and PTI constraints at mortgage origination, respectively. Constraint \eqref{eq:settle_debt} mandates that mortgage holders must repay at least a fraction $\upsilon_{1}$ of their mortgage principal. Lastly, constraint \eqref{eq:networth} equates the bequests left by deceased households to the sum of their financial and net housing wealth.

We denote the policy functions for consumption, savings, living with parents, becoming homeless, housing tenure, rental market choice, years living in the same unit, housing size, mortgage debt, mortgage origination, and bequests by $c_{j}(s)$, $a'_{j}(s)$, $b'_{j}(s)$, $v'_{j}(s)$, $o'_{j}(s)$, $f'_{j}(s)$, $\iota'_{j}(s)$, $h'_{j}(s)$, $m'_{j}(s)$, $n_{j}(s)$, and $q'_{j}(s)$, respectively.

\subsection{Rental company}\label{sec:rental_company}

We adopt the rental housing supply theory outlined in \citet{fallis-1985} and incorporate it within the rental housing supply framework developed by \citet{garrigaeetal-2009}. In our formulation, we designate $i=0$ to represent the controlled market and $i=1$ to represent the free market. Within each market $i$, for every rental housing unit, denoted by $x_i = 1$, the rental company obtains an average rent, $\overline{p}{r,i}$, while also bearing property taxes, $\tau{p}$, and depreciation costs, $\delta_{h}$.

In order to provide a total of $X'=\sum^{1}_{i=0} X'{i}$ rental housing, the rental company also faces management expenses amounting to $\frac{\theta_{1}}{\theta_{2}}X'^{\theta_{2}}$. $\theta_{1}$ signifies the average cost associated with managing the rental housing stock $X'$, while $\theta_{2}\geq 1$ delineates the price elasticity of the rental housing supply.

The rental company is obligated to provide a minimum quantity of controlled rental housing, which we denote as $\mbox{min}{\underline{X}_{0},X'}$, where $\underline{X}_{0} \geq 0$. If the total quantity of rental housing supplied, $X'$, is equal to or less than $\underline{X}_{0}$, the market falls entirely under control. In such instances, the rental company's marginal revenue equates to the average rent prevailing in the controlled market, denoted as $\overline{p}_{r,0}$. Conversely, if $X'$ exceeds $\underline{X}_{0}$, the rental company's marginal revenue becomes the average free market rent, denoted as $\overline{p}_{r,1}$. We use the term ``marginal rent" interchangeably with the rental company's marginal revenue. In scenarios in which $X'$ surpasses $\underline{X}_{0}$, $\underline{X}_{0}$ signifies the quantity of controlled units supplied, while $X' - \underline{X}_{0}$ represents the quantity of rental housing supplied in the free market.

The rental company purchases houses frictionlessly at the prevailing house price, $p$. In each period, it determines the quantity of rental housing to provide in the controlled market, represented by $X'_{0}$, and the quantity to supply in the free market, denoted as $X'_{1}$. The rental company is confronted with the following optimization problem:

\begin{align}\label{eq:rental_company_problem}
V(X)=\underset{X'_{0},X'_{1}} {\text{$\mbox{max} $}}\left\{\sum_{i=0}^{1} \overline{p}_{r,i}X'_{i}-\frac{\theta_{1}}{\theta_{2}}X'^{\theta_{2}}-p X'+p(1-\delta_{h}-\tau_{p})X+\frac{1}{1+r}V(X')\right\}.
\end{align}
Subject to
\begin{align}
& X'_{0}=\mbox{min}\{\underline{X}_{0},X'\} \label{eq:supply_restriction} \\
& X'_{1} \geq 0 \label{eq:supply_positive} \\
& X'=\sum^{1}_{i=0} X'_{i}. \label{eq:total_supply} 
\end{align}

\noindent Equation \eqref{eq:supply_restriction} stipulates that the quantity of rental housing supplied in the controlled market is the lesser of $\underline{X}_{0}$ and $X'$. Equation \eqref{eq:supply_positive} ensures that the quantity supplied in the free market cannot be negative. Equation \eqref{eq:total_supply} denotes that the total quantity supplied is the sum of the quantities supplied in both markets. When $X'>\underline{X}_{0}$, the supply function is expressed as follows:

\begin{align}\label{eq:rental_supply}
X'^{\star}=\left[\frac{1}{\theta_{1}}\left(\overline{p}_{r,1}-p\left(\frac{r+\delta_{h}+\tau_{p}}{1+r}\right)\right)\right]^{\frac{1}{\theta_{2}-1}}.
\end{align}

\noindent In equilibrium, when $X'^{\star} \geq \underline{X}_{0}$, it follows that $X'^{\star}_{0}=\underline{X}_{0}$ and $X'^{\star}_{1}=X'^{\star}-\underline{X}_{0}$. However, if $X'^{\star} < \underline{X}_{0}$ or if rent control laws are applicable to all rental units, then in equation \eqref{eq:rental_supply}, $\overline{p}_{r,0}$ replaces $\overline{p}_{r,1}$, and $X'^{\star}=X_{0}$ and $X'^{\star}_{1}=0$.

\noindent The partial derivatives of $X'^{\star}$ with respect to $p$ and $\overline{p}_{r,1}$ reveal the optimal quantity supplied:

\begin{itemize}
    \item Decreases with the housing price and increases with rent if $\theta_{2} > 1$, as is the case in our calibration.
    \item Remains independent of the housing price and rent if $\theta_{2} = 1$; in this scenario, rental supply is infinitely elastic.
\end{itemize}

One might contend that over time, the proportion of the controlled rental market naturally diminishes. This is due to the depreciation of existing rent-controlled units, which may be supplanted by new constructions not bound by rent control regulations. However, to sustain a consistent share of the rent-controlled market in a stationary equilibrium, governments can adjust raise the threshold year of construction beyond which units are exempt from rent control. This approach mirrors the actions taken by the Ontario government over the years, and is similar to measures implemented in Quebec, where newly constructed units are granted a five-year exemption from rent control regulations.

\subsection{Housing construction company}\label{sec:construction_company}

A competitive construction firm engages in the construction of new housing \citep{rotberg-2022}. It earns revenues, $pI$, from its production of new housing, denoted by $I$. However, the firm also bears costs, which are encapsulated by the following technology:

\begin{equation*}
TC(I)=\frac{\epsilon_{1}}{\epsilon_{2}}I^{\epsilon_{2}}\label{eq:total_cost_structure},
\end{equation*}

\noindent where $\epsilon_{1}$ dictates the average cost of housing investment, while $\epsilon_{2}$ governs the price elasticity of housing investment. The construction company makes decisions regarding investment to optimize its profits:

\begin{equation*}
\label{eq:construction_company_profits}
\underset{I} {\text{$\mbox{max}$}}{\left\{p I-TC(I)\right\}}.
\end{equation*}

\noindent The solution yields the following relationship between housing investment and the price of housing:

\begin{equation}
\label{eq:construction_company_first_order_condition}
I^{\star}=\epsilon_{1} p^{\epsilon_{2}-1}.
\end{equation}

\subsection{Aggregation}\label{subsection:aggregation}

As we outlined in the general description of the model, we incorporate both federal and local governmental entities. The federal government levies income and consumption taxes while also administering pensions to households. On the other hand, local governments collect property taxes and offer services to residents, including shelter for the homeless. Given that rent control policies are typically set at the local level, our focus remains on ensuring budgetary balance at this lower tier of government:

\begin{align}\label{eq:gbc}
  G + \sum_{j=1}^J \int_\mathcal{S} \theta_{b} b'_{j}(s)\mathrm{d}\Psi(s)  = \tau_{p} p H.
\end{align}

\noindent We require that government expenditures on services for residents and homeless support be matched by revenues from property taxes. When calibrating the model to align with the current tax code in Toronto, we determine the level of consumption of public goods, which we denote as $G$, to achieve a balanced budget. Throughout our experiments, we maintain the value of $G$ constant and adjust $\tau_{p}$, the property tax rate, to ensure budget balance.

$S=A\times V\times B\times O \times F \times I \times Z \times H\times M\times E$ denotes a typical subset of the state space $\mathcal{S}$; and $A$, $V$, $B$, $O$, $F$, $I$, $Z$, $H$, $M$, $X$ $E$ denote typical subsets of $\mathcal{A}$, $\mathcal{V}$, $\mathcal{B}$, $\mathcal{O}$, $\mathcal{F}$, $\mathcal{I}$, $\mathcal{Z}$, $\mathcal{H}$, $\mathcal{M}$, $\mathcal{X}$, and $\mathcal{E}$, respectively.

Aggregate market housing demand must equal aggregate housing supply:

\begin{equation}
\label{eq:housing_clearing_condition}
  \sum_{j=1}^J \int_\mathcal{S} (1-b'_{j}(s))(1-v'_{j}(s)) h'_{j}(s) \mathrm{d}\Psi(s) = H',
\end{equation}

\noindent and the housing stock law of motion is: 

\begin{equation}\label{eq:housing_stock_lom}
H'=I+(1-\delta_h)H.
\end{equation}


\noindent Equation \eqref{eq:housing_stock_lom} represents a standard law of motion for housing, stating that the next period's housing stock equals the housing investment plus the depreciated stock of housing. In each rental market, the aggregate rental housing demand must match the aggregate rental housing supply:

\begin{equation}
\label{eq:rental_housing_clearing_condition}
  \sum_{j=1}^J \int_\mathcal{S} (1-b'_{j}(s))(1-v'_{j}(s)) (1-o'_{j}(s))\indicator{f'_{j}(s)=i}h_{j}(s) \mathrm{d}\Psi(s) = X'_{i} \quad i \in \{0,1\}.
\end{equation}

The distributions of households evolve according to the laws of motion,
\begin{align}\label{eq:lom2}
\Psi_{1}(\mathrm{S})=\int_{\mathcal{E} \times \matchcal{X}} \indicator{(0,0,0,0,0,0,0,\underline{h}_{r},0,X,E)\in \mathrm{S}}\mathrm{d} \hat{F}(X)\tilde{F}(E),
\end{align}

\begin{align}
\Psi_{j+1}(S) &=\int_{\mathcal{S}} \hat{F}(X)\hat{F}(X)F(E,e|b'_{j}(s))G(z)M_{j}(s)\mathrm{d}\Psi_{j}(s),\ j=1,\ldots,J-1\label{eq:lom1}
\end{align}

\noindent where,

\begin{align}
    M_{j}(s) = \indicator{[a'_{j}(s)\in A] \land [v'_{j}(s)\in V] \land [b'_{j}(s)\in B] \land [o'_{j}(s)\in O] \land [f'_{j}(s)\in F] \land [\iota'_{j}(s)\in I] \land [h'_{j}(s) \in H] \land [m'_{j}(s) \in M]}.\notag
\end{align}

Finally, total bequests left by decedents must equal to the total wealth endowments:
\begin{equation}
\label{lom:agg_bequests}
      \int_{\mathcal{S}} q'_{J}(s) \mathrm{d}\Psi_{J}(s) =  \sum_{j=1}^J \int_{\mathcal{S}} \vartheta (\chi,e)\ \mathrm{d}\Psi_j(s).
\end{equation}

\noindent Similar to the methodology employed by \citet{kkk-2021}, we select the parameter $\vartheta$ to ensure this condition holds across all our experiments. Consequently, variations in $\vartheta$ between equilibria are contingent upon proportional changes in aggregate bequests.

\subsection{Equilibrium}
An equilibrium in the model is a sequence of aggregate prices and quantities, $\left\{H',I,\{X'_{i},\overline{p}_{r,i}\}^{1}_{i=0},p,q\right\}$, a sequence of value and policy functions, $\biggl\{V_{j}(\cdot),c_{j}(\cdot),a'_{j}(\cdot),v'_{j}(\cdot),b'_{j}(\cdot),o'_{j}(\cdot),f'_{j}(\cdot)\iota'_{j}(\cdot),h'_{j}(\cdot),m'_{j}(\cdot),n_{j}(\cdot),q'_{j}(\cdot)\biggl\}_{j=1}^J$, and a sequence of measures, $\left\{\Psi_{j}(\cdot)\right\}_{j=1}^J$, that:
\begin{enumerate}
\item Solve the household's dynamic problem (\ref{eq:vf}) subject to constraints (\ref{eq:bc})--(\ref{eq:networth});
\item Solve the rental company's problem (\ref{eq:rental_company_problem}) subject to (\ref{eq:supply_restriction})--(\ref{eq:total_supply});
\item Satisfy the construction company's first-order condition, (\ref{eq:construction_company_first_order_condition});
\item Satisfy the government's budget constraint, (\ref{eq:gbc});
\item Clear the housing market,
(\ref{eq:housing_clearing_condition}) subject to the housing stock law of motion (\ref{eq:housing_stock_lom});
\item Clear both rental housing markets,
(\ref{eq:rental_housing_clearing_condition});
\item Satisfy the laws of motion for the distribution of households, (\ref{eq:lom1})--(\ref{eq:lom2});
\item And satisfy the law of motion for bequests, (\ref{lom:agg_bequests}).
\end{enumerate}

\section{Calibration and model performance}\label{sec:calibration}

Our calibration involves two steps. Initially, we calibrate specific parameters using data from Toronto and estimates from the literature. Then, we calibrate the values of the remaining parameters to align the model's generated statistics with the data from Toronto. Table \ref{tab:calibration_external} details the parameter values we set externally along with their sources. Similarly, Table \ref{tab:calibration_internal} outlines the parameter values we set internally alongside their targets. In Section \ref{appendix:data} of the appendix, we elaborate on how we derived certain statistics using various datasets. Lastly, in Section \ref{subsection:non_targeted_statistics}, we demonstrate the model's performance across several dimensions not initially targeted.

\subsection{Assigned parameter values}

\subsubsection{Demographics, preferences, and income}
A model period spans 1 year. Households live from age 26 to age 85, and retire at age 66, which implies a retirement age of $J_{R}=41$ and a lifespan of $J=60$. Following \citet{rs-2023}, we set the maximum age at which households can live with their parents, $J_{p}$, to 9, thereby prohibiting households aged 35 and above from living with their parents. We set the households' relative risk aversion at $\sigma=2$. We established The equivalent household size for each age group, $\{\xi_{j}\}_{j=1}^{J}$, using 2016 Census data on the average household size for each age cohort. We scale household sizes by their square roots, following standard practice in the literature. We determine the discount factor $\beta$ and the utility weight on housing services $\gamma$ in the second stage of our calibration.

We adopt the labor productivity process outlined in \citet{gkkoc}. We set the standard deviation of the fixed effect to $0.128$, and the stochastic labor productivity shocks follow an AR(1) process in logarithmic terms,

\begin{equation*}
  \log e' = \rho_{e} \log e + \epsilon_{e},\quad \epsilon_{e}\overset{\mathrm{i.i.d.}} \sim N(0,\sigma_{e}).
\end{equation*}

\noindent We set the persistence parameter, $\rho_{e}$, to 0.9, following \citet{gkkoc}. In the second stage, we determine the standard deviation of the stochastic component of labor productivity, denoted as $\sigma_{e}$. Newborns obtain their stochastic labor productivities from the unconditional probability distribution associated with the ergodic distribution, which is endogenously generated by our productivity process. Consistent with \citet{az-2024}, for households who opt for or are compelled to be homeless, we mechanically increase the probability of remaining in their stochastic labor productivity state, $e$, by 20\% compared to what is implied by our aforementioned stochastic labor productivity process.

Based on the 2016 Census data, we establish the deterministic lifecycle component of labor market productivity as follows:

\begin{equation*}
\zeta_j = 0.923+0.09\times j-1.61\times 10^{-3}\times j^{2}, \mbox{ for } j\leq J_{R}
\end{equation*}

\noindent We derive this formula by regressing the average household income for each age, ranging from 26 to 65, on age and age squared. We determine pension benefits, denoted as ${y_{R}(\chi,e)}_{\{\chi \in \mathcal{X} \times e\in \mathcal{E}\}}$, based on the replacement rate of gross individual earnings, as outlined in \citet{OECD-pension}.

Surviving households' endowments are determined proportionally to their labor income, expressed as $\vartheta(\chi,e):=\vartheta \chi e$. We select the parameter $\vartheta$, governing surviving households' wealth endowments, to fulfill equation \eqref{lom:agg_bequests}.

\subsubsection{Taxes and interest rates}

We set the taxable portion of capital income, $\nu$, to be 50\%, mirroring the provision in Canada's tax code. Utilizing household-level data from the 2016 Census for Toronto, we estimate $\psi=0.9275$ in equation \eqref{eq:income_tax}. We set the consumption tax rate, denoted as $\tau_{c}$, to 13\% to align with the Harmonized Sales Tax (HST) rate in Ontario. Furthermore, we establish the interest rate, $r$, at 2.77\%, as estimated for Canada by \citet{bruneau-etal-2018}.

\subsubsection{Housing}

We set our rent shocks $\left\{z\right\}_{\mathbb{Z}}$ to $\left[-0.1227,-0.0607,0,0.0624,0.1284,0.2589\right]$, along with corresponding probabilities $G(z)$ given by $\left[0.16,0.12,0.44,0.08,0.16,0.04\right]$. We derive these values from 1992–2016 data sourced from the RMS, reflecting average rents in both controlled and free markets. Our analysis indicates a higher risk associated with substantial rent increases compared to the benefits of significant rent decreases.

To determine the rate at which real rents decrease with tenancy duration in the controlled market, denoted as $\kappa_{1}$, we refer to Census data, setting it at 1.365\%. Similarly, we set $\kappa_{2}$, representing the annual deterioration rate of controlled units, to 0.55\%, using the estimate from \citet{murray_et_al-1991} for Los Angeles.

We establish the per-household cost borne by the city to support the homeless, denoted as $\theta_{b}$, at 82.3\% of the average household income. We determine this parameter by dividing the City of Toronto's total spending on homelessness supports by the number of homeless individuals in the city in 2021.\footnote{See \url{https://www.homelesshub.ca/community-profile/toronto} and \url{https://www.toronto.ca/legdocs/mmis/2021/ex/bgrd/backgroundfile-163804.pdf}.}.

We set the property tax rate, $\tau_p$, to 0.44\% using information reported by Altus Group and 1992-2016 Property Sale and Assessment Data (PSAD). We set the housing depreciation rate to 0.6\% as estimated by \citet{rosenthal-2014}. Based on estimates by \citet{gm-2003}, we set buying costs, $\tau_{b}$, to 2\% of the purchase price, and selling costs, $\tau_{s}$, to 6\% of the sale price. We choose $\epsilon_{2}=1.89$ so that the price elasticity of housing supply is equal to $0.89$, as estimated by \citet{piaxao-2021} for Toronto. The level of the construction company's cost function, $\epsilon_1$, is a free parameter that we back out from the market clearing condition \eqref{eq:housing_clearing_condition}. Lastly, we set the weighted average rent in both rental markets, $\overline{p}_{r}$, to 5.24\% of the equilibrium house price, which we computed using 2016 data from the Canadian Real Estate Association (CREA).

We set the property tax rate, denoted as $\tau_p$, to 0.44\%, drawing upon information provided by Altus Group and 1992-2016 Property Sale and Assessment Data (PSAD). We set the housing depreciation rate, denoted as $\delta_{h}$ to 0.6\% based on estimates by \citet{rosenthal-2014}. Following \citet{gm-2003}, we determine buying costs, $\tau_{b}$, as 2\% of the purchase price, and selling costs, $\tau_{s}$, as 6\% of the sale price.

To achieve a price elasticity of housing supply equal to 0.89, as estimated by \citet{piaxao-2021} for Toronto, we select $\epsilon_{2}$ to be 1.89. Additionally, the construction company's cost function level, denoted as $\epsilon_1$, is a free parameter that we derive from the market clearing condition \eqref{eq:housing_clearing_condition}.

Lastly, we set the weighted average rent in both rental markets, denoted as $\overline{p}_{r}$, to 5.24\% of the equilibrium house price. We base our computation on 2016 data we obtained from the Canadian Real Estate Association (CREA).

\subsubsection{Mortgages}

In accordance with banking regulations in Canada, we set the loan-to-value (LTV) constraint, denoted as $\lambda_{m,1}(ph)$, as follows:

\begin{equation*}
\lambda_{m,1}(ph) = \begin{cases}
0.8 &\text{if $ph\geq \$1M$}\\
0.9 &\text{if $\$500K\leq ph< \$1M$}\\
0.95 &\text{if $ph < \$500K$}.
\end{cases}
\end{equation*}

\noindent Mortgage holders with a model-equivalent house value above \$1M can borrow only up to 80\% of its value; mortgage holders with a model-equivalent house value between \$500,000--\$1M can borrow up to 90\% of the house value; and mortgage holders with a house value below \$500,000 can get a loan of up to 95\% of the house value. We set $\lambda_{m,2}=0.39$ so that the maximum allowable PTI is 39\% as is standard in the mortgage industry in Canada.\footnote{See \url{https://www.cmhc-schl.gc.ca/consumers/home-buying/calculators/debt-service-calculator}.}

For mortgage holders with a model-equivalent house value exceeding \$1M, we restrict borrowing to 80\% of its value. Those with a model-equivalent house value ranging between \$500,000 and \$1M can borrow up to 90\% of the house value. Mortgage holders with a house value below \$500,000 are eligible for a loan of up to 95\% of the house value. We set $\lambda_{m,2}=0.39$ to align with the standard maximum allowable PTI of 39\% in the mortgage industry in Canada.\footnote{See \url{https://www.cmhc-schl.gc.ca/consumers/home-buying/calculators/debt-service-calculator}.}

In mortgage contracts, we begin by setting the fixed mortgage origination cost, denoted as $\omega_{1}$, to 0.68\% of the average household income in Toronto. We determine this based on data we source from the 2016 Census and the Canada Mortgage and Housing Corporation's (CMHC) website, which includes expenses such as appraisal fees (approximately \$250) and lawyer fees (approximately \$500).\footnote{See \url{https://www.cmhc-schl.gc.ca/en/professionals/industry-innovation-and-leadership/industry-expertise/resources-for-mortgage-professionals/how-much-will-my-home-really-cost}}

We set the prepayment penalty, denoted as $\omega_{2}$, to 3 months of interest, a standard practice in Canadian mortgage contracts. Following the methodology of \citet{rs-2023}, we choose $\upsilon_{1}$ and $\upsilon_{2}$. We set $\upsilon_{1}$, which represents the average fraction of mortgages paid each year, to 1.7\% to mirror a typical 30-year amortization period prevalent in the Canadian mortgage industry. We establish the threshold for the prepayment penalty, $\upsilon_{2}$, at $1.15 \times \upsilon_{1}$ plus 15\% of the total debt and interest payments, amounting to 18.8\%. This threshold is aligned with standard mortgage contracts in Canada, allowing prepayments of 15\% of the principal and regular payments without penalty.

Lastly, we set the interest rate premium on mortgage debt, denoted as $\Delta_{m}$, to 0.19\%. This figure represents the average spread between the 5-year variable rate on mortgages and the return on a 10-year government bond spanning from 2002 to 2022.\footnote{The mortgage data we use is not publicly available. The 10-year Government of Canada bond yield can be found at \url{https://www.bankofcanada.ca/rates/interest-rates/canadian-bonds/}}.

\subsection{Internally calibrated parameters}

After assigning the parameter values mentioned earlier, we proceed to calibrate another twelve parameters, as we outline in Table \ref{tab:calibration_internal}. These variables encompass the discount factor, $\beta$; utility weight on housing services consumption, $\gamma$; strength of bequest motive, $\mu$; average income taxation level, $\tau_{\ell}$; variability in labor income's stochastic component, $\sigma_{e}$; convexity of rental housing management, $\theta_{2}$; marginal managerial cost-to-rent ratio determined by $\theta_{1}$; initial rent paid by a tenant entering a unit under tenancy control, $p_{r,0}$; minimum rental unit size, $\underline{h}_{r}$; minimum housing unit size for purchase, $\underline{h}_{o}$; housing services for households living with parents, $\underline{h}_{p}$; and housing services for homeless households, $\underline{h}_{h}$.\footnote{The construction cost scalar, $\epsilon_{1}$, is a free parameter we choose to equate housing supply to housing demand in the benchmark equilibrium. In counterfactual simulations, we maintain $\epsilon_{1}$ at its benchmark level.}

We collectively calibrate all these parameters to align with twelve key metrics in Toronto data:

\begin{enumerate}
    \item The ratio of aggregate mortgage debt to aggregate house value among homeowners, which stands at 21.5\% according to the 2019 Survey of Financial Security (SFS).
    \item The proportion of total renters' income dedicated to rent, which we estimate at 22.5\% using 2016 Census data.
    \item The ratio of aggregate bequests to aggregate net wealth, which we calculate at 2.2\%, utilizing our estimates derived from the 2019 SFS.
    \item The average labor income tax rate, recorded at 19.6\%, which we source from the 2016 Census.
    \item The Gini coefficient of labor income, which we also obtain from the 2016 Census.
    \item The property tax burden on renters in cities akin to Toronto, which we estimate at approximately 12\%. We ascertain this figure by averaging estimates for similar cities from multiple sources, including \citet{cy-1994}, \citet{susin-2002}, \citet{fack-2006}, and \citet{sy-2020}.
    \item The ratio of free market rent to house price, which we calculate at 5.82\%, based on our analyses utilizing data from the 2016 Census and 2016 CREA.
    \item The percentage of renters in the controlled rental market, which we compute at 71\% using the 2016 Census. 
    \item The percentage of renters spending over 50\% of their before-tax income on rent, standing at 41\% according to the 2016 Census.
    \item The homeownership rate, reported at 66.5\% in the 2016 Census.
    \item The percentage of 26-34 year old households living with their parents, which we estimate at 22.3\% based on the 2016 Census.
    \item The percentage of households experiencing homelessness, which is documented at 0.26\% according to the "Homeless Hub".\footnote{See \url{https://www.homelesshub.ca/community-profile/toronto}.}
\end{enumerate}

While we determine all parameters jointly, each parameter tends to exert a more significant influence on specific outcomes. Here is how:
\begin{enumerate}
    \item Discount factor, $\beta$: Primarily shapes homeowners' debt levels; a lower $\beta$ leads to higher homeowner debt levels as it influences their optimal borrowing behavior.
    \item Utility weight on housing, $\gamma$: Influences the proportion of renters' income spent on housing; higher $\gamma$ values indicate a greater desire for housing services among households.
    \item Bequest motive strength, $\mu$: Affects the ratio of bequests to net wealth; higher $\mu$ values result in relatively more bequests compared to overall wealth.
    \item Average income taxation level, $\tau_{\ell}$: Determines the average income tax rate; lower $\tau_{\ell}$ values correspond to higher tax burdens on households' incomes.
    \item Variability in labor income's stochastic component, $\sigma_{e}$: Shapes labor income inequality; higher $\sigma_{e}$ values lead to a greater Gini coefficient for labor income distribution.
    \item Marginal cost of managing rental housing, $\theta_{2}$: Dictates the elasticity of rental housing supply; higher $\theta_{2}$ values result in less elastic rental supply and lower property tax burdens for renters.
    \item Marginal managerial cost-to-rent ratio, $\theta_{1}$: Influences the free market rent-to-price ratio; higher ratios are associated with higher $\theta_{1}$ values.
    \item Initial rent for controlled units, $p_{r,0}$: Controls the proportion of renters in the free market; higher initial rents relative to the free market lead to a larger share of renters opting for the free market.
    \item Minimum rental unit size, $\underline{h}_{r}$: Governs the proportion of renters spending over 50\% of their before-tax income on housing; larger minimum sizes result in a higher percentage of renters facing high housing cost burdens.
    \item Minimum housing unit size for purchase, $\underline{h}_{o}$: Influences the homeownership rate; larger minimum sizes lead to fewer households opting for homeownership.
    \item Housing services for households living with their parents, $\underline{h}_{p}$: Impacts the percentage of 26-34 year old households living with parents; higher services lead to a higher proportion of households residing with parents.
    \item Housing services for homeless households, $\underline{h}_{h}$: Affects the homelessness rate; greater services result in a higher rate of homelessness.
\end{enumerate}

\subsection{Model performance}\label{subsection:non_targeted_statistics}

We assess the performance of our model by examining aggregate, lifecycle, and tenure distribution statistics. Table \ref{tab:aggregate_non_targreted_statistics} presents the model's performance on aggregate statistics not specifically targeted during calibration.

Our model indicates a slightly higher ratio of average owned-house value to average household income compared to the data, alongside a somewhat lower percentage of owners with positive mortgage debt. Notably, our available data does not facilitate direct estimation of the ratio of average rental unit value to average owner-occupied unit value in Toronto. Hence, we compare our model's estimate with that of \citet{chatterjee-eyigungor-2015} for the United States, revealing a slight discrepancy.

Although directly comparable estimates are absent in our dataset, we include two additional statistics for reference. Our model projects a premium of approximately 6.2\% for new tenants in the controlled market, with average free market rents exceeding average controlled rents by 3\%. Furthermore, our model suggests that renters in the free market earn around 2.7\% more than their counterparts in the controlled market, a figure lower than observed in the 2016 Census.

The wealth distribution depicted in our model exhibits less concentration compared to real-world data. This divergence arises primarily from the absence of rate-of-return heterogeneity, a factor that \citet{fagereng} highlight as pivotal for driving wealth concentration at the upper echelons of the distribution. Our model also lacks features such as heterogeneous rate of time preferences \citep{ks-1998} or income shocks akin to those observed in studies such as \citet{mj-shocks}, which could further elucidate our inability to replicate high levels of wealth concentration.

Figure \ref{fig:life_cycle} illustrates the lifecycle outcomes of the model in comparison to empirical data. In panel (a), we observe that the homeownership rate across the lifecycle in our model is lower than the data for both young and old age groups, while slightly higher for middle-aged households. Moving to panel (b), young households in our model exhibit higher LTV ratios compared to their counterparts in the data, whereas the opposite trend holds true for older households.

Figure \ref{fig:rent_to_income_cdf} displays the cumulative frequency distribution (CDF) of rent-to-income ratios. Our model closely mirrors the rent-to-income ratios observed in the data, albeit with a slight lower number of households with high rent-to-income ratios and a slight higher number of households with low rent-to-income ratios.

Table \ref{tab:tenure_by_income_quintile} presents households' tenure categorized by their income quintiles. The second column from the left corresponds to the lowest income quintile, while the first column from the right represents the highest income quintile. We disaggregate tenure into two categories: the percentage of renters in each quintile residing in controlled rental units and the percentage residing in free market units. Panel (a) displays these statistics based on the data, while panel (b) showcases the calibrated model's performance regarding the same metrics.

Although the calibrated model does not perfectly replicate the data, it captures significant trends: notably, households in lower income quintiles are more inclined to rent in the controlled market compared to those in higher income quintiles.

\begin{comment}
Conferences in order of time:
1. Midwest Economics Association (MEA) in March 2023: https://mea.sites.grinnell.edu/ 
2. Midwest Macro Group in May 2023: https://krannert.purdue.edu/faculty/mcrucini/midwest-macroeconomics-group.php
3. Society for Economic Dynamics (SED) in June 2023: https://www.economicdynamics.org/
4. Canadian Macro Study Group (CMGS) in November 2023: https://canadianmacro.ca/
5. American Economic Association (AEA) in January 2024:
https://www.aeaweb.org/conference/
\end{comment}

\section{Experiments}\label{section:experiments}

\paragraph{Experiments' overview} Using our baseline calibrated model, we conduct four main experiments. In our first experiment we compare the benchmark of our baseline model, which features tenancy control alongside a free market, with the same model but with unit control instead of tenancy control. In the second experiment, we compare the baseline's benchmark (i.e., tenancy control alongside a free market) with the long-run equilibrium without any rent control, which we refer to as full decontrol. In the third experiment, we compare the baseline's benchmark with the long-run equilibrium in which all rental housing units are subject to tenancy control, which we refer to as full tenancy control. In the last experiment, we search for the degree of tenancy control that maximizes newborns' welfare in the long-run. In all our experiments, we keep the tax code fixed and balance the government's budget by adjusting the property tax rate. 

Using our baseline calibrated model, we conduct four experiments.

\begin{itemize}
    \item In our first experiment, we compare the benchmark of our baseline model, which incorporates tenancy control alongside a free market, with the long-run equilibrium of the same model featuring unit control instead of tenancy control.
    \item In the second experiment, we contrast the baseline benchmark with the long-run equilibrium in which there is no rent control, which we call full decontrol.
    \item For the third experiment, we compare the baseline benchmark with the long-run equilibrium in which all rental housing units are subject to tenancy control, which we refer to as full tenancy control.
    \item In the last experiment, we search for the degree of tenancy control that maximizes newborns' welfare in the long run.
\end{itemize}

Across all experiments, we keep the tax code fixed and balance the government's budget by adjusting the property tax rate.

\paragraph{Welfare measure} Our primary welfare measure is the expected ex-ante utility of newborns \citep{ralws-1971,conesa-et-al,gkkoc}. This measure is akin to providing all newborns with a fixed proportional non-housing consumption transfer. Formally, the non-housing consumption-equivalent welfare effect resulting from a policy change is expressed as:

\begin{equation}
  CE_{t} = \left[\frac{\sum_{j=1}^{k}\int_{\mathcal{S}}V^{\dagger}_{j}(s)\mathrm{d}\Psi^{\dagger}_{j}(s)-\sum_{j=1}^{k}\int_{\mathcal{S}}\omega_{B}(q'^{\star}_{J}(s))\mathrm{d}\Psi^{\dagger}_{j}(s)}{\sum_{j=1}^{k}\int_{\mathcal{S}}V^{\star}_{j}(s)\mathrm{d}\Psi^{\star}_{j}(s)-\sum_{j=1}^{k}\int_{\mathcal{S}}\omega_{B}(q'^{\star}_{j}(s))\mathrm{d}\Psi^{\star}_{j}(s)}\right]^{\frac{1}{(1-\gamma)(1-\sigma)}}-1,
\end{equation}

\noindent where index daggers denote counterfactual equilibrium objects, while stars denote objects in the baseline's benchmark. 

Additionally, we gauge the approval rate of the new policy change by quantifying the proportion of households experiencing an increase in their welfare under the new policy:

\begin{equation}
A = \sum_{j=1}^{k}\int_{\mathcal{S}} \indicator{V^{\dagger}_{j}(s)>V^{\star}_{j}(s)}\ \mathrm{d}\Psi^{\star}_{j}(s).
\end{equation}

\noindent For our primary welfare and approval measures, we set $k=1$, representing the welfare changes and approval rates of newborns. Additionally, we calculate average welfare changes and approval rates across all households by setting $k=J$. Since we compute approval rates based on the distribution in the baseline's benchmark, we can further assess the approval among homeowners and renters in the benchmark.


\paragraph{Measures of rent} Our analysis encompasses several measures of rent that we outline below:

\begin{itemize}
    \item \underline{Average free market rent}: The equilibrium average rent prevailing in the free market.
    \item \underline{Average controlled rent}: The equilibrium average rent observed in the controlled market.
    \item \underline{Rent of new tenants}: The equilibrium initial rent paid by new tenants upon entering a new rental unit. In the controlled market, new tenants face higher rents compared to long-standing tenants under rent control. In the free market, new tenants anticipate paying the same average rent as existing tenants in the free market.
    \item \underline{Average rent}: The weighted average of equilibrium rents across both the controlled and free rental markets.
\end{itemize}

\subsection{Tenancy control vs. unit control}

In our first experiment, we compare the baseline's benchmark with the long-run equilibrium featuring unit control. Before delving into the specifics of the model with unit control, we outline how we configure the unit control exercise to ensure comparability with the baseline's benchmark.

Firstly, we set the exogenously determined controlled rent of a unit in the unit control exercise to match the average controlled rent observed in our baseline's benchmark. Secondly, we align the proportion of the rental market under control in the unit control exercise with that of the baseline's benchmark. Lastly, to maintain consistency with the baseline while ensuring a constant real rent that is independent of tenancy duration in the unit control exercise, we set the rent decrease in the unit control exercise to mirror the unit's deterioration rate.

\subsubsection{Formal description of model with unit control}

We formalize the unit control model as follows: Each period, households ``wake-up" to encounter a random shock, which we denote as $\eta \in {0,1}$. The realization of $\eta$ yields a binary outcome: with a probability of $\pi$, $\eta$ equals 1; otherwise, it equals 0. Importantly, $\pi$ remains independent across time and previous realizations of $\eta$.

Households drawing $\eta = 1$ have the option to enter the controlled rental market in that period. Once in the controlled market, they retain the choice to stay in their current unit irrespective of subsequent realizations of $\eta$. However, upon vacating their unit, re-entry into the controlled market necessitates another realization of $\eta$ equal to 1. New tenants entering the controlled market initially pay a rent equivalent to the baseline average controlled rent, which stands approximately 2.9\% lower than benchmark free market rents and about 8.5\% lower than the initial rent paid by new tenants in the controlled market in the benchmark.

For each year a tenant remains in the same unit, their rent decreases by $\kappa_{1}=\kappa_{2}$, mirroring the rate at which their unit deteriorates due to landlord neglect. Consequently, although their nominal rent declines during their tenancy, their real rent remains constant, resulting in reduced housing expenses.

Conversely, households drawing $\eta = 0$ are ineligible to rent in the controlled market but can opt for the free market rental. Additionally, regardless of the $\eta$ realization, households have the choice to purchase a house, reside with parents if $j<J_{p}$, or face homelessness.

To maintain the same proportion of the rental market under control as in the baseline's benchmark, we choose the probability of drawing $\eta = 1$, denoted as $\pi$. Beyond these adjustments, the model remains consistent with the baseline's other assumptions.

\subsubsection{Aggregate results}

We present the aggregate findings of our initial experiment in the second column from the left in Table \ref{tab:aggregate_effects}. Given that we mechanically set the average controlled rent to match the average controlled rent in the benchmark, there would be no change in the average controlled rent. However, the allocation of controlled units under unit control would favor high-income households due to chance, which would lead to a notable 16.3\% increase in the homelessness rate. This shift is illustrated in panel (c) of Table \ref{tab:tenure_by_income_quintile}, which indicate a diminished proportion of low-income households renting in the controlled market compared to the benchmark.

Furthermore, the immediate rent discounts in the controlled market resulting from unit control, which take 10.5 years to materialize in the benchmark with tenancy control, would spur an increase in total rental demand at the intensive margin among remaining renters. Consequently, this surge would raise free market rents by 2.4\%, thereby decreasing the relative cost of homeownership, raising it by 0.15 percentage points and marginally decreasing overall rental demand at the extensive margin. The uptick in homeownership rates, coupled with the overall rise in total rental demand, would drive house prices up by 0.4\%. We depict a graphical representation of these dynamics in Figure \ref{fig:diagram_unit_control}.

Despite the fewer controlled units allocated to low-income households and the 2.4\% increase in free market rents, switching from tenancy control to unit control would only marginally decrease newborns' welfare by 0.06\%. Nevertheless, these losses would be widespread, benefiting less than 1\% of newborns. The welfare impact would be modest, primarily because unit control would promptly offer substantial rent discounts to renters fortunate enough to access the controlled market.

In a partial equilibrium simulation where we maintain rental prices at their benchmark levels while adjusting other parameters to their values under the long-run equilibrium with unit control, we observe a slight welfare improvement. Surprisingly, over 75\% of newborns would fare better under this scenario. This suggests that the immediate rent discounts facilitated by unit control would outweigh the drawbacks associated with the less favorable allocation of controlled units. However, this also implies that the general equilibrium effects of unit control on rental demand and free market rental prices would be significant enough to counterbalance these gains.

\subsubsection{Distributional results}

In the second row of Table \ref{tab:distributional_impact}, we provide a detailed breakdown of welfare changes for newborns across different income brackets. We segment the income distribution into five quintiles. As the table demonstrates, welfare losses would be evenly distributed across all income quintiles. Low-income households would predominantly lose due to the allocation of fewer controlled units to them under unit control. Additionally, unlucky low-income households would encounter higher and volatile rents in the free market, along with an increased risk of homelessness. These factors would overshadow the benefits stemming from immediate rent discounts accessible to fortunate low-income households in the controlled market. 

Conversely, high-income households would primarily experience losses due to higher house prices and increased free market rents. Given their tendency to swiftly transition to homeownership, the advantages of immediate rent discounts and a greater allocation of controlled units would be outweighed by these aforementioned drawbacks. Mid-income households would face a combination of these challenges, resulting in overall welfare losses.

\subsection{Benchmark vs. full decontrol}

\subsubsection{Aggregate results} The third column of Table \ref{tab:aggregate_effects} delineates changes in aggregate variables in the long-run equilibrium under full decontrol compared to the benchmark. Full decontrol would exert dual effects on the housing market. On one hand, it would diminish the incentive to remain in the same rental unit for an extended period, prompting more renters to transition to homeownership. This extensive margin effect would consequently raise homeownership, total housing demand, and house prices. Conversely, the elimination of the controlled market premium in the full decontrol equilibrium would lead to lower initial rents for remaining renters. Consequently, this would raise total rental demand among existing renters at the intensive margin, surpassing the impact of the extensive margin, and resulting in a 0.1\% increase in free market rents. We provide a visual representation of these dynamics in Figure \ref{fig:diagram_full_decontrol}.

Moreover, in the absence of the controlled market option in our full decontrol equilibrium, more households would be exposed to rental price volatility and higher average rents. This exposure would culminate in a substantial 32.2\% surge in the homelessness rate.

Decontrolling would adversely impact all newborns, reducing their welfare by 0.31\%. Additionally, it would lead to a decrease in welfare across the entire population by 0.14\%, with only 35.2\% of the population benefiting. These welfare losses would stem from two primary factors. Firstly, the removal of the controlled market would eliminate a valuable option that relatively efficiently allocated controlled units to low-income households. Secondly, it would result in increased house prices and free market rents, adversely affecting both prospective homeowners and renters who previously rented in the free market.

\subsubsection{Distributional results} 

In the third row of Table \ref{tab:distributional_impact}, we detail welfare changes across income quintiles. Under full decontrol, the reduction in newborns' welfare would be fairly uniform across the income distribution. Low-income households would primarily lose due to the elimination of the controlled market option, which is advantageous for them. Conversely, high-income households would experience losses primarily because of the higher free market rents and more expensive house prices, rendering homeownership less accessible. Mid-income households would face a combination of these factors, contributing to their welfare losses.

\subsection{Benchmark vs. full tenancy control}\label{subsection:full_control}

\subsubsection{Aggregate results}\label{sec:aggregate_results_full_tenancy_control}

In our third experiment, we analyze the long-term equilibrium under full tenancy control compared to the benchmark. The fourth column from the left of Table \ref{tab:aggregate_effects} outlines the changes in aggregate variables. Full tenancy control would compel more high-income households, along with some low-income households, who would otherwise opt for the free market, to enter the controlled market. This reduced exposure to volatility in free market rents would marginally decrease the homelessness rate.

Additionally, it would lead to more high-income households prolonging their tenure in the rental market compared to the benchmark. This elongation would result in a decrease in the homeownership rate by 0.26 percentage points and an increase in the demand for rental housing at the extensive margin.

Conversely, renters who had previously chosen the free market would now face higher initial rents in the controlled market, thereby reducing their rental demand at the intensive margin. The decline in rental demand stemming from the intensive margin would outweigh the increase in rental demand from the extensive margin, resulting in a relative decline in total rental demand compared to the benchmark by 0.6\%.

However, given that high-income households tend to transition to homeownership more swiftly than low-income households, forcing more high-income households into the controlled market would reduce the average duration of rental tenure and subsequently cause the average controlled rent to rise by 0.7\%. This adjustment would maintain initial rents roughly at the same level.

Lastly, as both rental demand and homeownership rates would decrease, housing prices would also experience a decline of 0.3\%. Figure \ref{fig:diagram_full_control} provides a qualitative diagram to visually illustrate these mechanisms.

We also find a slight decline in the welfare of newborns, equal to 0.02\%, although the average welfare across all households would slightly increase, by 0.16\%. The decrease in newborns' welfare would primarily stem from the removal of the free market option, which would reduce available rental choices for households. Typically, the presence of a free market option benefits mainly higher-income households, but occasionally some low-income households also prefer renting in the free market to avoid premium costs. Consequently, less than a quarter of all newborns would benefit from this policy change.

Conversely, the average welfare across all households would slightly rise because our calculation of average welfare includes households that have already been compelled into the controlled market and have accumulated some longer-term benefits of rent control. However, these benefits would extend only to approximately 27\% of the population.

\subsubsection{Distributional results}

In the fourth row of Table \ref{tab:distributional_impact}, we present a breakdown of the welfare gains and losses of newborns based on their position in the income distribution. Welfare losses are evident across all income quintiles. This would be attributed to the removal of the free market, a rental option that proves beneficial for numerous households, as we detailed in Section \ref{sec:aggregate_results_full_tenancy_control}.

\subsection{Optimal degree of tenancy control}\label{subsection:optimal_control}

We have demonstrated that, in our model, both complete decontrol and total tenancy control would result in overall welfare losses. This prompts the next logical inquiry: what level of tenancy control maximizes welfare? To answer this, we conduct a comprehensive grid search. We systematically vary the proportion of the market subject to tenancy control, and assess the welfare outcomes at each point within the corresponding long-run stationary equilibrium.

In our model, we identify the welfare-maximizing degree of tenancy control to be 65\%, which closely aligns with our benchmark in which approximately 70\% of the rental market falls under tenancy control. This level of control would yield only a marginal 0.01\% increase in welfare compared to the baseline benchmark. The optimal degree of tenancy control strikes a delicate balance between two competing objectives.

The first objective is to maintain a controlled market with affordable premiums, primarily benefiting low-income households by easing their housing costs and facilitating access to this segment of the market. As the proportion of the controlled market decreases, competition for controlled units intensifies, resulting in higher premiums that burden low-income households.

The second objective aims to keep rents in the free market affordable for households who prefer renting outside the controlled market, which predominantly includes higher-income households, but also some low-income households who opt for this option. Achieving this goal involves reducing the size of the controlled market. However, beyond a certain threshold, an expansion of the controlled market compels more households to rent within it. Consequently, these households become entrenched in renting for longer periods due to the benefits accrued from rent control with each additional rental term. This phenomenon leads to an overall increase in rental demand and subsequently drives up free market rents. In our model, we find that these trade-offs are effectively balanced when 65\% of the rental market is subject to tenancy control.

For completeness, we include the aggregate results of this last experiment in the fifth column of Table \ref{tab:aggregate_effects} and the distributional outcomes in the fifth row of Table \ref{tab:distributional_impact}. However, we refrain from delving into these results, as the disparities in relation to the benchmark are inconsequential.

Given the proximity of the welfare-maximizing degree of tenancy control to the benchmark level, which we have already compared to unit control, we opt not to delve further into its disparities relative to unit control. However, we did replicate the exercise with unit control and observed outcomes akin to those of our initial experiment.

\section{Conclusion}\label{sec:conclusion}

We constructed a general equilibrium model encompassing both tenancy control and unit control --- two prevalent rent control policies --- to analyze their respective impacts on welfare and equilibrium prices. Leveraging household-level Census data and household-level SFS data from Toronto, where approximately 71\% of rental units fall under tenancy control while the remainder remain uncontrolled, we calibrated our model to reflect the local housing market dynamics.

Utilizing the calibrated model, we conducted four distinct experiments to evaluate the impacts of various rent control policies. In the first experiment, we compared the benchmark against the long-run equilibrium under unit control, maintaining an identical degree of control and average controlled rents. The second experiment compared the long-run equilibrium under full tenancy control with the benchmark. In the third experiment, we contrasted the benchmark with the long-run equilibrium under full rent decontrol. Our final experiment involved pinpointing the welfare-maximizing degree of tenancy control in the long run.

Across these experiments, we observed notable disparities. Unit control would allocate fewer controlled rental units to low-income households compared to tenancy control. Furthermore, unit control would lead to an increase in total rental demand, consequently elevating free market rents, and thus marginally diminishing welfare and exacerbating homelessness relative to the benchmark. Both full decontrol and full control would result in a reduction in welfare across the entire income spectrum relative to the benchmark, primarily due to the elimination of viable rental options for households. Finally, determining the welfare-maximizing degree of tenancy control aimed to strike a balance between reducing free market rents and maintaining affordable initial rents in the controlled market. In our model, we identified this equilibrium at 65\% of the rental market under tenancy control, a degree closely aligned with our benchmark.

\textcolor{red}{Something new which we have yet to tell the reader.}

\textcolor{red}{1. Maybe something about that tenancy control is endogenous and thus does not results in large rent discounts on average. Thus, under tenancy control the housing stock may be better maintained in the long-run (something we did not capture with our model). With unit control government can set very low rents.}

\textcolor{red}{2. Maybe something about how in NYC there is some evidence that some landlords are acting as if they are under tenancy control --- from a welfare point of view it may be good.}


\bibliographystyle{aer}
\bibliography{reference_files/bib}

\vspace{1cm}
    
\begin{table}[H]
    \centering
    \caption{Externally assigned parameter values}
    \vspace{-1.25cm}
    \label{tab:calibration_external}
    \input{tables/calibration_external}
    \vspace{-0.5cm}
    \begin{fignote2}{1.11\textwidth}
    \textit{Notes:} Table \ref{tab:calibration_external} presents the values of the externally calibrated parameters. Part (i) delineates the external calibration for non-housing parameters, while Part (ii) outlines the external calibration of housing parameters. Additionally, Part (iii) provides the external calibration of mortgage-related parameters.
    \end{fignote2}
\end{table}

\begin{table}[H]
    \centering
    \caption{Internally assigned parameter values}
    \vspace{-1.25cm}
    \label{tab:calibration_internal}
    \input{tables/calibration_internal}
    \vspace{-0.5cm}
    \begin{fignote2}{1.08\textwidth}
    \textit{Notes:} Table \ref{tab:calibration_internal} displays the values of the internally calibrated parameters. Our benchmark is closely aligned with all of our targets. From the 2019 SFS we source data on the percentage of homeowners with mortgages and the aggregate bequests to aggregate net wealth ratio. From the 2016 Census, we obtain data for Toronto CMA regarding the total rent to total income ratio, the average income tax rate, the labor income Gini coefficient, the percentage of renters in the controlled market, the percentage of renters spending over 30\% of their income on rent, the homeownership rate, and the percentage of 26-34-year-old households living with their parents. Additionally, using 2016 Census and 2016 CREA data for Toronto CMA, we calculate the free market rent to housing price ratio. We source the tax incidence of renters from \citet{cy-1994}, \citet{susin-2002}, \citet{fack-2006}, and \citet{sy-2020}, and obtain the homelessness rate from the ``Homelessness Hub".
    \end{fignote2}
\end{table}

\begin{table}[H]
    \centering
    \caption{Aggregate non-targeted statistics}
    \vspace{-0.5cm}
    \label{tab:aggregate_non_targreted_statistics}  \input{tables/aggregate_non_targeted_statistics.tex}
    \vspace{-0.45cm}
    \begin{fignote2}{0.9\textwidth}
      \textit{Notes:} Table \ref{tab:aggregate_non_targreted_statistics} presents the model's performance on aggregate and distributional non-targeted statistics compared to the data. The second column from the left displays the model's performance, the third column includes estimates from the literature or the data, and the last column provides the source. "--" indicates that the statistics is unavailable.
    \end{fignote2}
\end{table}

\begin{table}[H]
    \centering
    \caption{Rental market choice by income quintile (all households)}
    \vspace{-0.5cm}
    \label{tab:tenure_by_income_quintile}    
    \input{tables/tenure_by_income_quintile_everyone.tex}
    \vspace{-0.45cm}
    \begin{fignote2}{0.66\textwidth}
      \textit{Notes:} Table \ref{tab:tenure_by_income_quintile} illustrates rental market tenure by income quintile for all households that we computed using Toronto CMA data from the 2016 Census. Panel (a) presents these statistics in the data, panel (b) showcases the baseline model's performance, and panel (c) displays outcomes under unit control. The second column from the left represents the lowest income quintile, while the first column from the right signifies the highest quintile. The first row in each panel indicates the percentage of households in each quintile renting in the controlled market. For instance, in panel (a), 28.88\% of all households in the second quintile rent in the controlled market. The second row in each panel indicates the percentage of households in each quintile renting in the free market.  
    \end{fignote2}
\end{table}

\begin{table}[H]
    \centering
    \caption{Aggregate effects generated by changes in rent control policy}
    \vspace{-0.5cm}
    \label{tab:aggregate_effects}
    \input{tables/aggregate_effects.tex}
    \vspace{-0.45cm}
    \begin{fignote2}{1.1\textwidth}
      \textit{Notes:}  In Table \ref{tab:aggregate_effects}, the second column contrasts the benchmark with the long-run equilibrium in which we enforce, via unit control, a price ceiling on controlled units that is identical to the average controlled rent in our benchmark. Additionally, in that exercise, we set the proportion of the rental market under control to equal to that of our benchmark. The third column juxtaposes the benchmark with the long-run equilibrium under full decontrol. The fourth column compares the benchmark with the long-run equilibrium under full tenancy control. The last column compares the benchmark with the welfare-maximizing degree of tenancy control in the long-run. "—" indicates incompatible statistics. We assess welfare in non-housing consumption-equivalent terms.
    \end{fignote2}
\end{table}

\begin{table}[H]
    \centering
    \caption{Distributional impacts of rent control policy changes on newborns' welfare, broken down by income quintile ($\%\Delta$ in welfare).}
    \vspace{-0.5cm}
    \label{tab:distributional_impact}
    \input{tables/distributional_impacts.tex}
    \vspace{-0.45cm}
    \begin{fignote2}{0.77\textwidth}
      \textit{Notes:} Table \ref{tab:distributional_impact} presents welfare changes for newborns, comparing unit control, full decontrol, and full tenancy control to the benchmark. We measure welfare in non-housing consumption-equivalent terms.
    \end{fignote2}
\end{table}

\begin{figure}[H]
  \centering
    \caption{Diagram illustrating rent trajectories under tenancy control based on rental market selection and duration in the same rental unit.}
    \vspace{-0.37cm}
    \label{fig:diagram_rent_path}
    \includegraphics[width=\textwidth]{figures/path_of_rent_diagram.pdf}
    \vspace{-1.3cm}
    \begin{fignote2}{1.0\textwidth}
      \textit{Notes:} The downward-sloping blue line depicts rental prices paid by tenants residing in controlled units, which are a function of their duration of occupancy in those units. Conversely, the horizontal red line represents the average rental price faced by tenants opting for the free market, with the dashed red line representing rent fluctuations they may face in free market rents. Notably, the average rental price in the free market remains constant regardless of tenancy duration. For clarity in this illustration, we abstract from the deterioration of units in the controlled market. Early during a tenancy, tenants in the free market pay lower average rents than those in controlled units, but later in the tenancy this reverses. Why is this the case? If initial rents in the controlled market ($p_{r,0}$) were equal to or lower than free market rents, then all renters would logically opt for controlled rentals. To entice some tenants towards the free market, those in controlled units must initially pay a premium to enjoy the safeguarding against rent fluctuations and lower long-term rental prices. Tenants anticipating short-term stays and willing to tolerate risk, opt for the free market to avoid this premium. Conversely, tenants planning extended stays and having low risk tolerance, are willing to pay the initial premium to capitalize on the advantages of rent control.      
    \end{fignote2}
\end{figure}

\begin{figure}[H]
  \centering
    \caption{Lifecycle statistics}
    \vspace{-0.2cm}
    \label{fig:life_cycle}
    \includegraphics[width=\textwidth]{figures/lifecycle.pdf}
    \vspace{-0.45cm}
    \begin{fignote2}{1\textwidth}
      \textit{Notes:} In Figure \ref{fig:life_cycle}, the two panels illustrate the model's performance on lifecycle statistics compared to the data. Panel (a) displays the average homeownership rate across all ages, while panel (b) presents the average loan-to-value (LTV) ratios for homeowners across 12 age groups, as the 2019 SFS provides household maintainer's age in 5-year ranges. Model outputs are depicted in orange, while data points are represented in blue. We source homeownership rates by age for Toronto CMA from the 2016 Census, and obtain data on the average LTV ratios for Ontario from the 2019 SFS.
    \end{fignote2}
\end{figure}

\begin{figure}[H]
  \centering
    \caption{Rent-to-income ratio cumulative frequency distribution}
    \vspace{-0.2cm}
    \label{fig:rent_to_income_cdf}
    \includegraphics[width=\textwidth]{figures/rent_income_cdf.pdf}
    \vspace{-0.45cm}
    \begin{fignote2}{1\textwidth}
      \textit{Notes:} In Figure \ref{fig:rent_to_income_cdf}, we depict the cumulative distribution function (CDF) of the rent-to-income ratio for renters. We represent the model's output in orange, and the data in blue. We calculate the rent-to-income CDF using data from the 2016 Census for Toronto CMA.
    \end{fignote2}
\end{figure}

\begin{figure}[H]
  \centering
    \caption{Qualitative diagram illustrating the impact of switching from the baseline's benchmark with tenancy control to unit control}
    \vspace{-0.36cm}
    \label{fig:diagram_unit_control}
    \includegraphics[width=\textwidth]{figures/unit_control_diagram.pdf}
    \vspace{-0.6cm}
    \begin{fignote2}{1.0\textwidth}
      \textit{Notes:} In Figure \ref{fig:diagram_unit_control}, we present a qualitative diagram illustrating the impact of switching from the benchmark to unit control with the same average controlled rent as in the benchmark. The switch to unit control would shift the rental supply curve leftward from the solid blue line, denoted as $\text{Rental supply}_{1}$, to the dashed blue line, denoted as $\text{Rental supply}_{2}$. This shift results from the increase in equilibrium house prices and property taxes (stemming from the rise in homelessness rates), which increases the user-cost of housing. This is evident in equation \eqref{eq:rental_supply}, where an increase in house prices or property taxes results in a decrease in the supply of rental housing. The switch to unit control would induce two effects additional effects: one at the intensive margin (pertaining to the decisions of existing renters) and one at the extensive margin (relating to the choice between owning and renting). Unit control would offer immediate rent discounts to renters in the controlled market, discounts that would take renters in the controlled market approximately 10.5 years to achieve in the benchmark. This would increase total rental demand at the intensive margin, shifting it from the solid red line, denoted as $\text{Rental demand}_{1}$, to the dotted red line, denoted as $\text{Rental demand}_{2}$. Combined with the shift in supply, this would lead to an increase in free market rents, thereby slightly raising homeownership and reducing total rental demand at the extensive margin, moving it from the dotted red line, $\text{Rental demand}_{2}$, to the dashed red line, $\text{Rental demand}_{3}$. In the benchmark equilibrium, we denote the quantity as $X^{1}$, the average free market rent as $\overline{p}^{1}_{r,1}$, average controlled rent as $\overline{p}^{1}_{r,0}$, and average rent as $\overline{p}^{1}_{r}$. Under unit control equilibrium, the quantity increases to $X^{2}>X^{1}$, with free market rent rising to $\overline{p}^{2}_{r,1}>\overline{p}^{1}_{r,1}$.
      
    \end{fignote2}
\end{figure}

\begin{figure}[H]
  \centering
    \caption{Qualitative diagram illustrating the impact of full rent decontrol}
    \vspace{-0.36cm}
    \label{fig:diagram_full_decontrol}
    \includegraphics[width=\textwidth]{figures/full_decontrol_diagram.pdf}
    \vspace{-0.6cm}
    \begin{fignote2}{1.0\textwidth}
      \textit{Notes:} In Figure \ref{fig:diagram_full_decontrol}, we present a qualitative diagram illustrating the impact of full decontrol compared to the benchmark. Similar to the analysis provided in Figure \ref{fig:diagram_unit_control}, the supply curve would shift leftward due to the same factors elaborated in the notes of that figure. As in the case of unit control, full decontrol would also exert two effects on rental demand: one at the extensive margin and one at the intensive margin. Decontrol would diminish the incentive to rent, leading to increased homeownership and consequently lower total rental demand. This shift would be represented by the movement of the rental demand curve to the left, from the solid red line, denoted as $\text{Rental demand}_{1}$ to the dashed red line, denoted as $\text{Rental demand}_{2}$. Simultaneously, decontrol would eliminate the rent premium for all new tenants, prompting each remaining renter to demand larger rental units. This would cause a counteracting shift in rental demand to the right from the dashed red line, $\text{Rental demand}_{2}$, to the dotted red line, denoted as $\text{Rental demand}_{3}$, representing the final location of rental demand. This latter shift would be driven by changes at the intensive margin, which would exert a stronger influence than the extensive margin. In the benchmark equilibrium, we denote quantity as $X^{1}$, average free market rent as $\overline{p}^{1}_{r,1}$, average controlled rent as $\overline{p}^{1}_{r,0}$, and average rent as $\overline{p}^{1}_{r}$. Under full rent decontrol equilibrium, the quantity increases to $X^{2}>X^{1}$, with free market rent rising to $\overline{p}^{2}_{r,1}>\overline{p}^{1}_{r,1}$, and the average rent, which would be equivalent to the free market rent in the new equilibrium, $\overline{p}^{2}_{r,1}>\overline{p}^{1}_{r}$.
      
    \end{fignote2}
\end{figure}

\begin{figure}[H]
  \centering
    \caption{Qualitative diagram illustrating the impact of full tenancy contro}
    \vspace{-0.36cm}
    \label{fig:diagram_full_control}
    \includegraphics[width=\textwidth]{figures/full_tenancy_control_diagram.pdf}
    \vspace{-0.8cm}
    \begin{fignote2}{1.0\textwidth}
      \textit{Notes:} Figure \ref{fig:diagram_full_control} provides a qualitative diagram depicting the impact of full tenancy control relative to the benchmark. In contrast to full decontrol, full tenancy control would shift the rental supply curve rightward from the solid blue line, denoted as $\text{Rental supply}_{1}$, to the dashed blue line, denoted as $\text{Rental supply}_{2}$. This shift would occur because equilibrium house prices would decrease and the property tax rate, $\tau_p$, would also slightly decline due to a reduction in the homelessness rate. Equation \eqref{eq:rental_supply} illustrates this, indicating that rental housing supply increases as house prices and property taxes decline. Full tenancy control would induce two additional effects on rental demand: an extensive margin effect and an intensive margin effect. Encouraged by the availability of controlled units, more households would opt for renting, and the rental demand curve would shift rightward from the solid red line, $\text{Rental demand}_{1}$, to the dashed red line, $\text{Rental demand}_{2}$, due to changes at the extensive margin. Subsequently, because all new tenants would be compelled to pay a premium in the controlled market, renters who previously occupied larger units in the free market would now demand smaller units. This intensive margin effect would shift rental demand leftward from the dashed red line, $\text{Rental demand}_{2}$, to the dotted red line, $\text{Rental demand}_{3}$, representing the final rental demand position. Although the intensive margin effect would outweigh the extensive margin effect, decreasing total rental demand relative to the benchmark, the total rental demand for controlled units would surpass that in the benchmark. Consequently, this would raise the average controlled rent. In the benchmark equilibrium, we denote the quantity as $X^{1}$, free market rent as $\overline{p}^{1}_{r,1}$, average controlled rent as $\overline{p}^{1}_{r,0}$, and average rent as $\overline{p}^{1}_{r}$. Under full tenancy control equilibrium, the quantity would increase from $X^{1}$ to $X^{2}$, the average controlled rent would also increase to $\overline{p}^{1}_{r,0}$ to $\overline{p}^{2}_{r,0}$, and average rent, which would be the average-controlled rent in the new equilibrium would rise to $\overline{p}^{2}_{r,0}<\overline{p}^{1}_{r}$. In the benchmark, the quantity supplied is determined by the free market rent, $\overline{p}^{1}_{r,1}$, which represents the marginal rent. Under full tenancy control, the quantity supplied would be determined by the average controlled rent, denoted as $\overline{p}^{2}_{r,1}$. Since $\overline{p}^{2}_{r,1}<\overline{p}^{1}_{r,1}$, the quantity supplied under full tenancy control, $X^{2}$, would be lower than the quantity supplied in the benchmark, $X^{1}$.      
    \end{fignote2}
\end{figure}
    
\newpage

\appendix

\maketitle

\counterwithin{figure}{section}
\counterwithin{table}{section}
\counterwithin{equation}{section}

\part*{Appendix (For Online Publication)}

Most of our assigned parameter values and calibration targets are taken directly from the literature. However, there are still many that we constructed ourselves from several data sources. In Section \ref{appendix:data} we describe some of the more complex ones, breaking them down to rental market variables, income-related variables, and other variables. All the variables are computed for Toronto, CMA except variables computed using the 2019 SFS, where we use data for Ontario because the lowest geography available in the publically available dataset is province.

\section{Data}\label{appendix:data}

\subsection{Rental markets}

\paragraph{Rent shocks in free market}

To establish the rent shocks in the free market and their associated probabilities, we leverage data from the 1992-2016 Rental Market Survey (RMS), which is not publicly available. For each year, we calculate the average rent in both the controlled and free markets. Next, we run the time-series data for each market through an HP-filter, applying a penalty parameter of 1,600, to generate long-term trend time-series for each market.

Subsequently, for each year and market, we compute the percentage deviation of actual rent from its respective long-term trend. Then, we compute the standard error of deviations from trend in the controlled market relative to the average rent in that market, utilizing it to adjust each free-market deviation accordingly. This adjustment ensures comparability between markets.

Following this, we determine the occurrences in which the percentage deviation from trend in the free market falls more than two standard errors below the trend. From these occurrences, we compute the average percentage deviation, which serves as the value of $z$. We repeat this process for deviations falling within different ranges of standard errors below and above the trend.

Ultimately, this methodology allows us to ascertain both the probability of occurrence for $z$ and its corresponding value across various deviations from trend.

\paragraph{Percentage of renters in the controlled market}

To compute the percentage of renters residing in controlled units, we rely on data from the 2016 Census. We utilize the variable BUILT to distinguish between controlled and free-market units, subsequently dividing the total number of controlled units by the total number of rental units across both markets. This analysis reveals that 71\% of renters lease properties in the controlled market, indicating that 29\% of renters opt for units in the free market. We denote the fraction of renters in the controlled market as $D$.

\paragraph{Real rent decrease, $\kappa_{1}$, for each year remaining in controlled units}

In our model, we need to calibrate the rate at which real controlled rents decrease  with existing renters, which we denote in our paper by $\kappa_{1}$. To compute this, we first need to know the average number of years a renter remains in the same rental. 

Using 2016-2021 RMS data, we compute an average renter turnover rate of 12.25\%. Utilizing this value, we compute the average number of years a renter stays in the same unit using the following formula:

\begin{equation}\label{eq:avg_tenure_duration}
    Y=\sum_{i=1}^{50}i(1-(1-\tau_{r})^{i}).
\end{equation}

\noindent $Y$ represents the average number of years renters stay in the same unit, $\tau_{r}$ denotes the turnover rate, and $i$ signifies the number of years.

To grasp the mechanics of equation \eqref{eq:avg_tenure_duration}, imagine an initial mass of renters equal to $1$. After one year, a fraction of $1-\tau_{r}$ remains in their units, implying that a fraction of $\tau_{r}$ resided in their unit for precisely one year. After two years, a fraction of $(1-\tau_{r})^2$ continues to reside in their unit, suggesting that a fraction of $1-(1-\tau_{r})^2$ dwelled in their unit for exactly two years, and so forth. We truncate the calculation at 50 years, as the majority of renters do not remain in the same unit for such an extended duration. By multiplying each fraction of renters living precisely $i$ years in the same unit by the number of years they resided in their unit, we obtain the expected number of years renters live in the same unit.

Next, using the estimated fraction of renters in the controlled market ($D$) and our approximation for the average number of years a renter stays in the same unit across both the controlled and free market ($Y$), we can derive the average number of years a renter in the controlled market remains in the same unit, $Y_{c}$, as follows:

\begin{equation*}
    Y_{c}=\frac{Y-(1-D)Y_{f}}{D}.
\end{equation*}

\noindent $Y_{f}$ denotes the average number of years a renter stays in a free market unit. To obtain $Y_{c}$, we leverage the estimate in \citet{cjq-2023} that renters in the controlled market remain in the same unit twice as long as renters in the free market, which implies that $Y_{f}=\frac{Y_{c}}{2}$. Thus, we can solve for $Y_{c}$.

Once we have solved for the value of $Y_{c}$, we incorporate two additional data points from the Census. Firstly, we consider the average rent in the controlled market, utilizing the variables SHELCO and BUILT. Secondly, we estimate the average rent in the controlled market among renters who moved in the past year, using the variables SHELCO, BUILT, and Mob1. From our analysis, we estimate that renters who recently moved into their controlled unit pay approximately $13.84\%$ more than the average rent across the entire controlled market, denoted as $Q$. Given $Q$, we estimate $\kappa_{1} = 0.01365$ using the following equation:

\begin{equation}\label{eq:computing_annual_decrease}
    \kappa_{1}=((1+Q)^{\frac{1}{Y_{c}}}-1).
\end{equation}

\paragraph{Rent to house price ratio}

To calculate the rent-to-house price ratio, we first compute the average annual rents across all rental units in the 2016 RMS data. Then, we divide this figure by the average apartment price in 2016, as reported in CREA data. Please note that CREA data is an internal data source and not publicly available. This computation results in a ratio of $0.0524$.

\paragraph{Rent to house price ratio in free market}

To compute the rent-to-house price ratio in the free market, we utilize three estimated figures. Firstly, we have our computed rent-to-house price ratio across both markets, which is $0.0524$ as discussed earlier. Secondly, we use our estimate of the share of the controlled market in the total rental market, which we determined to be $D=29\%$. Thirdly, we incorporate an estimate from \citet{zhang-gou} regarding how much lower the average rent in the controlled market is relative to the average rent in the free market, all else being equal. We denote this difference as $T$. We formulate the rent-to-price ratio as follows:

\begin{equation}\label{eq:free_market_rent_to_price}
    {\left(\frac{\overline{p}_{r}}{p}\right)}=D \left(\frac{p^{0}_{r}}{p}\right)+(1-D) \left(\frac{p^{1}_{r}}{p}\right).
\end{equation}

\noindent $\left(\frac{\overline{p}_{r}}{p}\right)$ represent the rent-to-average price ratio. Given that $\frac{p^{0}_{r}}{p}=(1-T)\frac{p^{1}_{r}}{p}$, we can substitute this into equation \eqref{eq:free_market_rent_to_price} and solve for $\frac{p^{1}_{r}}{p}=0.0582$.

\paragraph{Tenure distribution by income quintile}

We utilize 2016 Census data to determine the percentage of households in each income quintile who rent in the free market and who rent in the controlled market. To categorize households into income quintiles, we utilize the total income variable (TOTINC). To distinguish renters, we employ a binary variable indicating whether a household is a renter or an owner (TENUR), and to differentiate renters in the controlled and free markets, we utilize the categorical variable representing the year of dwelling construction (BUILT).

\subsection{Income and taxes}

\paragraph{Increase in real average income by age}

To derive the deterministic age-related growth component of labor income, $\zeta_{j}$, we leverage 2016 Census data. Firstly, we calculate the average market income (variable name: EMPIN) for households whose household head's age ranged between 25 and 64, respectively (variable name: AGEGRP). Subsequently, regress income on age and age squared, as we outlined in Section \ref{sec:calibration}.

\paragraph{Average income tax rate}

We calculated the average income tax rate based on data from the 2016 Census. This involved dividing the total taxes paid (variables: TOTINC - TOTINC\_AT) by the total income (variable name: TOTINC). The resulting average income tax rate is 19.5\%.

\paragraph{Estimating progressivity of labor income taxes, $\psi$}

In our model, we define taxes by equation \eqref{eq:taxable_labor_income}. Rearranging this equation, we obtain:

\begin{equation}\label{eq:taxable_labor_income_regression}
  \mbox{log}(y^{x}_{j}(e,a)-y^{\tau}_{j}(e,a))=\mbox{log}(\tau_{\ell})+\psi \mbox{log}(y^{x}_{j}(e,a)).
\end{equation}

\noindent The left-hand side variable represents after-tax taxable household income, while the right-hand side variable denotes taxable household income. The intercept is $\log(\tau_{\ell})$, and the coefficient on taxable household income is $\psi$, which determines the progressivity of taxes. Using 2016 Census data in equation \eqref{eq:taxable_labor_income_regression}, we obtain a point estimate for $\psi$ equal to 0.927. Since taxes in the 2016 Census were not categorized by income type, we utilize after-tax household income from all sources and before-tax household income from all sources. In this context, our dependent variable is after-tax household income (variable name: TOTINC\_AT), while the independent variable is before-tax household income (variable: TOTINC).

\subsection{Other statisitcs}

\paragraph{Property tax rate}

To calculate the property tax rate, we draw on two sources. Firstly, we consider the average property tax rate compiled by Altus Group, which they report as 0.66\%.\footnote{See \url{https://go.altusgroup.com/l/575253/2022-10-24/343cbn}.} However, property values assessed by municipalities in the City of Toronto often differ significantly from their actual sale prices. To address this, we analyze house-level 1992--2016 Property Sale and Assessment Data (PSAD). Although this dataset is not publicly available, it provides insights into the sale prices and assessment values of individual properties. By examining this data, we compute the average assessment value to sale price ratio, revealing that assessment values are, on average, 33\% lower than sale values. Consequently, this implies an effective property tax rate of 0.44\%.

\paragraph{Aggregate bequests to aggregate net wealth}

To compute the aggregate bequest to aggregate net wealth ratio, we use data from the 2019 Survey of Financial Security (SFS) and the Life Tables for Canada, which we sourced from the World Health Organization.\footnote{See \url{apps.who.int/gho/data/view.searo.60290?lang=en}.} We calculate the net wealth for economic families within each age cohort (recorded under the variable name: PWNETWPG) and then multiply this aggregate sum by the death rate of each age cohort. This process provides us with the total value of bequests for each age cohort. Summing up all bequests across all age cohorts, we divide this total by the sum of net wealth for all economic families. As a result, we obtain a ratio of 0.022, indicating that aggregate bequests constitute 2.2\% of aggregate net wealth. This estimate aligns closely with figures reported for the United States (see \citet{cagetti-denardi-2006}).

\end{document}
